\documentclass[11pt]{article}
\usepackage[T1]{fontenc}
\usepackage{amsmath,amssymb,amsthm,mathtools}
\usepackage{geometry,hyperref}
\geometry{margin=1in}
\setlength{\emergencystretch}{1em}
\hypersetup{colorlinks=true,linkcolor=blue,urlcolor=blue,citecolor=blue}
\newtheorem{theorem}{Theorem}
\newtheorem{proposition}[theorem]{Proposition}
\newtheorem{lemma}[theorem]{Lemma}
\newcommand{\D}{\mathbb D}
\newcommand{\C}{\mathbb C}
\newcommand{\E}{\mathbb E}
\newcommand{\Cat}{\operatorname{Cat}}
\newcommand{\capc}{\operatorname{cap}}
\title{Springer coefficients, boundary moments, polynomial zeros, and regular stars}
\hypersetup{pdftitle={Springer coefficients, boundary moments, polynomial zeros, and regular stars},pdfauthor={Cheng Li}}
\author{Cheng Li\\
	\footnotesize School of Mathematical Sciences, Fudan University,\\
	\footnotesize Shanghai 200433, China\\
	\footnotesize E-mail: \texttt{24300180118@m.fudan.edu.cn}}
\date{}
\begin{document}
	\maketitle
	\begin{abstract}
		For centered exterior univalent maps we prove the Springer
		inverse-coefficient bound in every order indexed by \(k\geq751\),
		with equality precisely for unit slit maps. The energy estimate
		uses an explicit multiplier that is linear up to one cutoff and
		constant thereafter. We determine the sharp radial boundary-moment
		maxima and give quantitative deficits for real exponents
		\mbox{\(4\leq p\leq8\)}. In the same range, the normalized derivative-zero
		moments of polynomials with connected lemniscates have a strict
		upper bound equal to their supremum across degrees. The regular
		\(m\)-arm star maximizes the polynomial factor constant among
		\(m\)-fold rotationally symmetric plane continua precisely when \(m\)
		is even. We give a quantitative even-order bound and classify
		equality, and construct a strict counterexample in every odd
		order \(m\geq3\).
	\end{abstract}
	
	\section{Introduction}
	The normalized exterior map links coefficient problems in univalent
	function theory with equilibrium measures of plane continua.
	Springer \cite{Springer} proposed that the modulus of the coefficient
	of order $2k-1$ of the inverse at infinity is at most $\Cat_{k-1}$.
	The area theorem gives the first order, and Springer proved the bound for $B_3$.
	Earlier low-order bounds are due to Kubota and Schober
	\cite{Kubota,Schober77}. Schober's corollary
	\cite[p.~413]{Schober} gives the sharp bounds through $B_{15}$,
	including the slit equality cases. With the sign convention below, Theorem~\ref{thm:large} treats
	every odd index from $1501$ onward. Its proof uses signed area and energy
	identities for the actual map, followed by a comparison near the two
	slit endpoints. In the energy estimate two coefficients of a single
	flux are chosen proportional to the same capped linear function.
	This reduces the multiplier comparison to a quadratic inequality.
	
	Baernstein, Laugesen, and Pritsker proved a sharp comparison for
	one-dimensional projections of equilibrium measure
	\cite[Theorem~2]{BLP}. Section~6 of their paper proves the second
	radial moment bound with its rotated-segment extremals and formulates
	a broader radial comparison as Conjecture~2. Its hypotheses are that
	the continuum has capacity one, contains zero, and has conformal
	centroid zero. Its tests are $C^1$ functions $\phi$ for which both
	$\phi$ and $\phi'$ are convex, applied to $\log|z|$. The test
	$\phi(x)=e^{px}$ gives the conjectured comparison for $p>0$ under
	those hypotheses. Corollary~6.3 proves the logarithmic comparison
	for arbitrary convex tests when the capacity-one continuum contains
	zero and is symmetric about it. We prove the bounds for $4\leq p\leq8$
	for every centered exterior map, with a quantitative deficit.
	The boundary argument extracts a second-moment deficit
	from the area theorem at the second moment. A finite weighted area
	inequality and low-order Grunsky estimates then give the fourth,
	sixth and eighth moments. Interpolation retains a uniform deficit
	throughout \([4,8]\). A comparison with the actual exterior map of
	a polynomial lemniscate transfers this bound to the zero moments.
	The comparison of the $p$th moments of polynomial zeros with those
	of their derivatives is classical for $p\geq1$ \cite[Theorem~4]{DBS1947}.
	Here the strict bound comes from excluding the slit as the exterior of
	a polynomial sublevel set. Scaled Chebyshev polynomials show that
	the bound is the supremum even after any fixed number of derivatives.
	
	The polynomial factor constant is governed by the equilibrium
	average of the logarithmic farthest-point distance
	\cite[Theorem~1.1]{PR2}. Pritsker and Ruscheweyh proved a regular-star
	bound under inclusions and conjectured that the inclusions could be
	dropped \cite[Theorem~2.2 and the discussion after Theorem~2.3]{PR2}.
	For connected compact sets we obtain the comparison in every even
	order and a counterexample in every odd order $m\geq3$. Their separate
	question about regular polygons concerns convex sets. The even star
	comparison uses the power map,
	the symmetric logarithmic moment comparison, and the second-moment
	deficit. In odd order, an explicit deformation gives a strictly
	increasing integral lower bound for the factor constant. The
	appendices give the exterior normalization and the analytic
	comparisons used in the verification.
	
	\section{Large-order inverse coefficients}\label{sec:large}
	Let $F(z)=z+b_1z^{-1}+b_2z^{-2}+\cdots$ be univalent on
	$\{|z|>1\}$, with the constant Laurent coefficient zero.  Write its
	inverse at infinity with the sign convention
	\[
	G(w)=w-\sum_{j\geq1}B_jw^{-j}.
	\]
	Set $\Cat_q=(q+1)^{-1}\binom{2q}{q}$.  A unit slit map means
	$F_a(z)=z+a/z$, $|a|=1$; equality of maps below is on the whole open
	exterior disk.
	
	\begin{theorem}[Large-order Springer statement]\label{thm:large}
		For every centered exterior univalent $F$ and every integer $k\geq751$,
		\[
		|B_{2k-1}(F)|\leq\Cat_{k-1}.
		\]
		For each fixed $k\geq751$, equality holds if and only if
		$F=F_a$ on $\{|z|>1\}$ for some $a\in\C$ with $|a|=1$.
	\end{theorem}
	
	We use $n=2k-1\geq1501$, $m=n-1\geq1500$.  An endpoint deficit is denoted
	$d\geq0$, with $t=nd$, $\Delta=md$, and
	$L=\log(2/d)=\log(2m/\Delta)$.  The zero-deficit case must be handled
	by geometric slit rigidity before the positive-deficit parameters are
	introduced.
	
	\subsection{Actual-map analytic bridge}\label{sec:bridge}
	Put $g(u)=\sum_{j\geq1}b_ju^j$, $H(z)=z+g(\bar z)$ on the unit disk,
	and $d\sigma=dA/\pi$.  Let $K$ be the actual omitted continuum.
	The following identities fix the sign and the normalization used by
	every subsequent saving calculation.  Put
	$\mathcal M_m=\int_{\D}|H|^m\,d\sigma$ and
	$\mathcal A_m=\int_K|w|^m\,d\sigma$.
	\begin{equation}\label{eq:signed-fill}
		B_{2k-1}=\int_{\D}H^mg'(\bar z)\,d\sigma,
		\qquad \mathcal M_m-\mathcal A_m
		=\int_{\D}|H|^m|g'(\bar z)|^2\,d\sigma.
	\end{equation}
	For $0<s<1$ put $F_s(z)=sF(z/s)$. Its unit-circle image is a
	smooth Jordan curve, with positive orientation around the bounded
	complement. The filling $H_s(z)=z+s g(s\bar z)$ has this boundary
	trace and Jacobian $1-|s^2 g'(s\bar z)|^2$. On the boundary the
	residue formula for the inverse gives
	\[
	B_n(F_s)=\frac{1}{2\pi i n}\oint_{|z|=1}H_s^n\,dz .
	\]
	Cauchy--Green gives the first formula in \eqref{eq:signed-fill},
	since $\partial_{\bar z}H_s=s^2g'(s\bar z)$. For the second, the form
	$|w|^m(x\,dy-y\,dx)/(m+2)$ has exterior derivative $|w|^m\,dx\,dy$.
	Apply Stokes' theorem both to its pullback by $H_s$ and to the
	bounded domain with the same oriented boundary.
	
	To take $s\uparrow1$, set $E_s=\C\setminus F(\{|z|>1/s\})$.
	These sets decrease to $K$, while the set omitted by $F_s$ is $sE_s$.
	For any fixed $s_0>0$ with $s_0<1$, $E_{s_0}$ is bounded. Thus
	\[
	\int_{K(F_s)} |w|^m\,d\sigma(w)
	=s^{m+2}\int_{E_s}|w|^m\,d\sigma(w)
	\longrightarrow \int_K|w|^m\,d\sigma(w).
	\]
	The area theorem gives $\sum_{j\geq1}j|b_j|^2\leq1$, so the
	orthogonality of disk monomials yields
	$s^2g'(s\bar z)\to g'(\bar z)$ in $L^2(d\sigma)$.
	Their squared moduli converge in $L^1$ by Cauchy--Schwarz.
	Also $H_s\to H$ pointwise and $|H_s|\leq2$: the normalized
	exterior omission bound gives $|w|\leq2$ on $K(F_s)$, and $H_s$
	is the Poisson extension of its boundary values. The omission bound
	itself follows by applying the area theorem to $\sqrt{F_s(z^2)-w}$,
	whose first negative coefficient is $-w/2$. These uniform bounds and
	dominated convergence pass the disk integrals to the limit. Finally
	$G_s(w)=sG(w/s)$ gives $B_n(F_s)=s^{n+1}B_n(F)$.
	
	Choose $|\beta|=1$ with $\beta^2B_{2k-1}=|B_{2k-1}|$
	(take $\beta=1$ if $B_{2k-1}=0$) and set
	$U=\operatorname{Im}(\beta H^k)$.  Expanding the gradient square
	and using \eqref{eq:signed-fill} yields the exact identity
	\begin{equation}\label{eq:signed-energy}
		\int_{\D}|H|^m\,d\sigma-|B_n|
		=\frac12\int_K|w|^m\,d\sigma
		+\frac1{2k^2}\int_{\D}|\nabla U|^2\,d\sigma.
	\end{equation}
	Lemma~\ref{lem:actual-phase-energy} proves
	$k^2\int U^2/|H|^2\,d\sigma\leq
	(7/3)\int|\nabla U|^2\,d\sigma$
	and gives the signed integrand used below.  Its proof uses a drift
	identity, one combined flux, and direct cutoff integration for the
	actual filling.  The choice of multiplier and its scalar estimate are
	given in Section~\ref{sec:energy-cert}.
	With $C=7/3$ set
	\[
	S(w)=\left(\frac{|w|}{2}\right)^m
	\left(1-\frac{\sin^2(\arg(\beta w^k))}{2C}\right),
	\quad S(0)=0.
	\]
	Discarding the nonnegative omitted-area term from
	\eqref{eq:signed-energy} then gives the signed comparison
	\begin{equation}\label{eq:signed-criterion}
		2^{-m}|B_n(F)|\leq\int_{\D}S(H)\,d\sigma.
	\end{equation}
	The one-sided reference moment is
	\begin{equation}\label{eq:reference-moment}
		M_m:=\int_{\D}(\operatorname{Re}z)_+^m\,d\sigma
		=\frac{\binom{m}{m/2}}{2^m(m+2)},
		\qquad 2^{m+1}M_m=\Cat_{k-1}.
	\end{equation}
	Consequently a candidate with $|B_n|\geq\Cat_{k-1}$ requires
	$\int S(H)\,d\sigma\geq2M_m$.  Every strict comparison below
	uses this dimensionless normalization.
	
	The radial trace $F^*$ of Lemma~\ref{lem:trace} takes values in $K$
	almost everywhere. The trace of $H$ is the same, since
	$H(se^{it})-F(s^{-1}e^{it})=(s-s^{-1})e^{it}$ for $0<s<1$.
	The bounded harmonic filling is the Poisson extension of this trace.
	The trace pushes normalized angular measure to the equilibrium measure
	of $K$, or harmonic measure at infinity. The normalization gives
	$K$ capacity one and conformal centroid zero. Thus the projection
	theorem \cite[Theorem~2]{BLP} puts
	$u(t)=\operatorname{Re}(\bar v F^*(e^{it}))$ below $2\cos t$ in convex
	order for each fixed $|v|=1$.
	
	\paragraph{Poisson comparison for fixed projections.}
	Write $d\mu=dt/(2\pi)$ and take convolution with respect to $\mu$.
	If a bounded boundary function $u$ is below an even, decreasing
	profile $q$ in convex order, then
	\begin{equation}\label{eq:poisson-convex-order}
		P_r*u\preceq_{\rm cx}P_r*q,\qquad 0<r<1 .
	\end{equation}
	For an angular set $E$ let $E^*$ be the centered arc of mass $\mu(E)$,
	and let $u^*$ be the even decreasing rearrangement of $u$.
	Circular rearrangement \cite{Circle} bounds $\int_E P_r*u\,d\mu$ by
	$\int u^*(P_r*\mathbf 1_{E^*})\,d\mu$. A common additive constant
	handles signed functions. The weight $P_r*\mathbf 1_{E^*}$ is even
	and decreasing on $[0,\pi]$: for $E^*=[-b,b]$ its derivative has
	the sign of $P_r(t+b)-P_r(t-b)$, which is nonpositive for
	$0\leq t,b\leq\pi$. Its level sets are centered arcs.
	For any arc mass $s$, the integral of $u^*$ on that arc is at
	most that of $q$. Indeed the upper integral of $u$ over sets of
	mass $s$ is the infimum of $as+\int(u-a)_+\,d\mu$ over
	$a\in\mathbb R$, and the same holds for $q$. Integrating over
	the levels of the weight gives
	$\int_E P_r*u\,d\mu\leq\int_{E^*}P_r*q\,d\mu$.
	The means are equal. Taking the supremum of
	$\int_E(P_r*u-a)\,d\mu$ over $E$ compares the positive-part
	tests at all $a$. These tests and equality of the means give
	convex order by piecewise linear approximation on the bounded range.
	
	With $q(t)=2\cos t$, its Poisson extension is $2r\cos t$. Thus
	\begin{equation}\label{eq:ordinary-convex}
		\E_t\phi(\operatorname{Re}(\bar v H(re^{it})))
		\leq\E_t\phi(2r\cos t)\quad(0<r<1)
	\end{equation}
	for every convex real $\phi$. The positive-part test at $2r$ and
	continuity in the angle give $\operatorname{Re}(\bar vH(z))\leq2|z|$.
	Using all unit directions yields $|H(z)|\leq2|z|$. The Poisson
	formula also gives $H(\D)\subset\operatorname{conv}K$.
	The two-point geometry is stronger than the one-point bound:
	\begin{equation}\label{eq:two-point}
		4|A-P|+|A+P|^2\leq16
		\qquad(A,P\in\operatorname{conv}K).
	\end{equation}
	For $A,P\in K$, the normalized odd maps
	$h_A(z)=\sqrt{F(z^2)-A}$ and $h_P$ are exterior univalent.
	When $P=A$, \eqref{eq:two-point} is the omission bound $|P|\leq2$.
	For $P\ne A$ put $q=(P-A)/4$ and $V=(h_A+h_P)/2$.
	The identity $h_A^2-h_P^2=4q$ makes $V$ nonvanishing. Moreover
	$h_A=V+q/V$ and $h_P=V-q/V$. Equal $V$-values force equal
	$h_A$-values, so $V$ too is exterior univalent. The area theorem and the
	parallelogram identity give
	$|q|^2\mathcal E(1/V)\leq A_V$, where
	$\mathcal E(f)=\sum_{j\geq1}j|[z^{-j}]f|^2$ and
	$A_V=\operatorname{area}(K_V)/\pi=1-\mathcal E(V)$.
	The first coefficient of $1/V$ is one, so $A_V\geq |q|^2>0$.
	By the Laurent area integral and change of variables under $V$,
	\[
	\mathcal E(1/V)=\frac1\pi\int_{|z|>1}|(1/V)'(z)|^2\,dA(z)
	=\frac1\pi\int_{\C\setminus K_V}|w|^{-4}\,dA(w).
	\]
	Let $D_V$ be the centered disk of area $\pi A_V$. The sets
	$D_V\setminus K_V$ and $K_V\setminus D_V$ have equal area, and the
	weight on the first is at least $A_V^{-2}$ and on the second at most
	$A_V^{-2}$. The last integral is therefore at least the normalized
	integral over $\C\setminus D_V$, which is $1/A_V$.
	The first negative coefficient of $V$ gives
	$A_V\leq1-|A+P|^2/16$.  These imply \eqref{eq:two-point}.
	Joint convexity extends it to the whole ordinary convex hull.
	If $P\in K$ has $|P|=2$, equality in the area theorem forces
	$F(z)=z+P^2/(4z)$, establishing the zero-deficit slit case.
	
	For a selected endpoint $P=R\xi$, $|\xi|=1$ and $R=2-d$, put
	$F_\xi(z)=\bar\xi F(\xi z)$ and $H_\xi(z)=\bar\xi H(\xi z)$.
	The rotation sends $P$ to $R$ and $K$ to $\bar\xi K$.
	The rotated real part of $w$ is $\operatorname{Re}(\bar\xi w)$.
	Rotating the argument and value preserves univalence and
	$F_\xi(z)=z+O(z^{-1})$. Its filling is
	$H_\xi(z)=z+\bar\xi g(\bar\xi\bar z)$.
	Let $A=R/2$, $E=1-A^2$ and, for $r>1$, set
	\[
	p_r=r-A/r,\quad
	S_r=\tfrac12\log\frac{r^2+1}{r^2-1}-r^{-2},\quad
	W=(p_r+\sqrt{E S_r})^2.
	\]
	The odd univalent root $\sqrt{F_\xi(z^2)-R}$ has first negative
	coefficient $-A$ and a tail with area budget $E$. For $w\in K$,
	each root of $\bar\xi w-R$ is omitted by this map. It therefore lies
	in the closed bounded region enclosed by the image of $|z|=r$.
	On that image the real part of $z-A/z$ has absolute value at most
	$p_r$, and Cauchy--Schwarz bounds the tail by $\sqrt{ES_r}$.
	The same bound $\sqrt W$ on the absolute real part holds in the
	bounded region. Writing $x+iy=\bar\xi w$ and squaring the root gives
	\begin{equation}\label{eq:whole-strip}
		y^2\leq4W(W+R-x),
		\qquad |w|^2\leq(x-2W)^2+4WR.
	\end{equation}
	The first sublevel region is convex, which justifies its extension
	from $K$ to $\operatorname{conv}K$.
	
	For a projected boundary function $u$ with $u\leq2-d$, $0<d<2$,
	let $h\in(0,\pi)$ satisfy $2\sin h/h=2-d$. Let $u_d$ equal
	$2\cos t$ outside $[-h,h]$ and its average $2-d$ inside.
	For any angular set of mass $s$, the $u$-integral is at most both
	$(2-d)s$ and the centered-arc integral of $2\cos t$. Their minimum
	is the centered-arc integral of $u_d$. For $s\leq h/\pi$ the arc lies
	inside the plateau. For $s\geq h/\pi$ it contains the entire plateau,
	whose integral is unchanged. Both boundary means are zero. The upper
	partial-mean criterion gives $u\preceq_{\rm cx}u_d$, so
	\eqref{eq:poisson-convex-order} gives the capped Poisson comparison.
	The direction stays fixed on each source circle.
	\subsection{The scalar cutoff}
	The following precise input is derived from an \emph{actual} contender
	in the boundary localization below:
	\begin{equation}\label{eq:candidate-input}
		\begin{gathered}
			|B_n(F)|\geq\Cat_{k-1},\quad d>0
			\quad\Longrightarrow\quad 0<t<1,\\
			J_n(t):=0.067nt^{5/2}e^{-0.61t}-0.285z-0.275z^2<0.58,\\
			z=0.34t\log(2n/t).
		\end{gathered}
	\end{equation}
	The implication holds for \emph{each} positive endpoint deficit,
	with the weak candidate inequality shown.  The boundary argument
	obtains it from the joint estimate
	$J_n(t_1)+J_n(t_2)<0.52$ and the following scalar floor.
	
	\paragraph{The scalar floor.}
	For $n\geq1501$ and $0<t<1$, put
	$L=\log(2n/t)$ and $z=0.34tL$.  The elementary envelope
	$(c+u)^3\leq c^3e^{u/2}$ for $c\geq6$, $u\geq0$ follows from
	$\log(1+u/c)\leq u/c$.  Apply it first with
	$c=\log(2n)$, $u=-\log t$, and then with
	$c=\log3002$, $u=\log(n/1501)$.  The resulting base check,
	using $\log3002<8.1$, gives
	\[
	70z^3\leq nt^2\sqrt t.
	\]
	Since $e^{-0.61t}>8/15$, the positive term of $J_n$
	strictly exceeds $(5/2)z^3$.  Finally
	\[
	\frac52z^3-\frac{11}{40}z^2-\frac{57}{200}z+\frac9{160}
	=(z-\tfrac14)^2(\tfrac52z+\tfrac9{10})
	+\tfrac3{40}z^2+\tfrac7{800}z\geq0.
	\]
	Thus $J_n(t)>-9/160$.  The joint bound makes each
	$J_n(t_i)<0.52+9/160<0.58$, including a candidate at equality.
	
	\subsection{Two actual boundary ends and the scalar necessity}
	This subsection proves the boundary input to the scalar cutoff
	by explicit inequalities.  Assume $F$ is not a unit slit,
	$k\geq751$, $n=2k-1$, and
	$|B_n(F)|\geq\Cat_{k-1}$.  Put
	$I_n=\E|2\cos t|^n$.  The inverse residue formula on an exterior
	circle, followed by bounded radial convergence, gives
	\begin{equation}\label{eq:boundary-candidate}
		nB_n=\E[e^{it}F^*(e^{it})^n],\qquad
		\E|F^*|^n\geq n\Cat_{k-1}=I_{n+1}/2
		\geq(1-0.502/n)I_n.
	\end{equation}
	For the last estimate, under the $|\cos t|^n$ weight the first two
	moments of $\sin^2t$ are $1/(n+2)$ and
	$3/((n+2)(n+4))$.  Apply
	$1-\sqrt{1-u}\leq(u+u^2)/2$ and check
	$(n+7)/(2(n+2)(n+4))\leq0.502/n$ for $n\geq1501$.
	Wallis' bound (or integration over $|t|\leq3/\sqrt n$) gives
	\begin{equation}\label{eq:In-bounds}
		0.785\frac{2^n}{\sqrt n}\leq I_n
		\leq\sqrt{2/\pi}\frac{2^n}{\sqrt n}.
	\end{equation}
	
	Choose $P_1\in K$ of maximal modulus, say
	$R_1=|P_1|=2-d_1$ and $\xi_1=P_1/R_1$.  The following short
	argument supplies both a stronger first radius and an actual second
	endpoint before any fine two-sided moment comparison.
	
	\begin{lemma}[An explicit opposite endpoint]\label{lem:imaginary-opposite}
		For every candidate above, $R_1>359/180$.  There is a point
		$w\in\operatorname{conv}K$ with
		$\operatorname{Re}(\bar\xi_1w)<-8/5$.  Consequently the maximum-modulus
		point $P_2$ in the opposite closed half of $\operatorname{conv}K$
		satisfies $|P_2|>8/5$ and belongs to $K$.  The ordinary-hull strip
		applies to both $P_1$ and $P_2$; this does not yet assert the later
		$nd_i<1$ deficit bounds.
	\end{lemma}
	\begin{proof}
		The actual residue and omitted-radius bound gives
		$n\Cat_{k-1}\leq R_1^n$.  Put $A=R_1/2$ and
		$c_k=\binom{2k}{k}/4^k=n\Cat_{k-1}/2^n$.
		Here $c_1=1/2$ and
		$c_{k+1}/c_k=(2k+1)/(2k+2)$; squaring the recurrence and using
		$(2k+1)^2\geq4k(k+1)$ proves $c_k\geq1/(2\sqrt{k})$.
		For $a=359/360$, the binomial theorem gives
		\[
		a^{-1501}=(1+1/359)^{1501}
		\geq\sum_{j=0}^{6}\binom{1501}{j}359^{-j}
		>55>2\sqrt{751}.
		\]
		Thus $a^{1501}<c_{751}$.  The ratio
		$(2k+1)/(2k+2)\geq1503/1504>a^2$ for $k\geq751$ propagates
		$a^{2k-1}<c_k$ by induction.  Since $A^n\geq c_k$, we obtain
		$A>a$, and hence $1-A^2<1/180$.
		
		Rotate by $\xi_1$, so $P_1=R_1>0$ on the real axis.  The actual odd
		square-root map $h$ obeys
		$h(z)^2=F_{\xi_1}(z^2)-R_1$,
		$b_1(h)=-A$, and its remaining odd Laurent coefficients have
		weighted square budget at most $1-A^2$.  Set
		$q=25/24$, $r=\sqrt q$, $M=r+A/r$, and
		$T=h(ir)-iM$.  The odd-mode Cauchy--Schwarz bound at this one point is
		\[
		|T|^2\leq(1-A^2)
		\left\{\frac12\log\frac{1+24/25}{1-24/25}-\frac{24}{25}\right\}
		=(1-A^2)(\log7-24/25)<\frac1{180}.
		\]
		The last logarithmic comparison follows from
		$e^{49/25}\geq\sum_{j=0}^{6}(49/25)^j/j!>7$.
		Thus $|T|<3/40$ and $M\leq49/24$.  The identities
		$(ir)^2=-q$ and $F_{\xi_1}(-q)=R_1+(iM+T)^2$ give, directly from
		the definition of the actual harmonic filling,
		\begin{equation}\label{eq:imaginary-root-witness}
			H_{\xi_1}(-1/q)=F_{\xi_1}(-q)+q-1/q
			=-\frac{1+A^2}{q}+2iMT+T^2.
		\end{equation}
		The exact rational bounds
		$(1+a^2)24/25>153/80$ and
		$2(49/24)(3/40)+(3/40)^2<5/16$ imply
		$\operatorname{Re}H_{\xi_1}(-1/q)<-8/5$.
		As $1/q<1$, the actual filling lies in the ordinary convex hull;
		rotating back gives the stated $w$.  The selected opposite-half
		maximizer has $|P_2|\geq|w|>8/5$.
		
		The root tail bound holds on the entire source circle $|z|=r$.
		The squared strip width therefore satisfies
		\[
		W\leq(r-A/r+3/40)^2<1/64,
		\qquad r-A/r=(q-A)/r<2/45.
		\]
		At the cutting line $x=\operatorname{Re}(\bar\xi_1v)=0$ the full-hull
		strip gives $|v|^2\leq4W(W+R_1)<1/4$.  Since $|P_2|>8/5$, it is
		inside the negative half-plane, where the strict convexity of
		$|\cdot|^2$ makes it an extreme point of $\operatorname{conv}K$.
		Compactness of $K$ puts that extreme point in $K$.  Applying the
		actual omitted-point strip to each $P_i$ proves the last assertion.
	\end{proof}
	
	\begin{lemma}[A Gaussian single-end exclusion]\label{lem:second-radius-195}
		For the same selected $P_2$ in Lemma~\ref{lem:imaginary-opposite},
		$R_2=|P_2|>39/20$.
	\end{lemma}
	\begin{proof}
		Put $c=39/20$, $q=39/40$ and $e=2-|w|$.  The two-point inequality
		\eqref{eq:two-point} holds for $w,P_1$ throughout the ordinary convex
		hull.  If $|w|\geq c$, then $0\leq e\leq1/20$ and
		$u=|w-P_1|$ satisfies $u(4-u)\leq16e$.
		The elementary quadratic alternative is
		$u\leq17e/4$ or $4-u\leq17e/4$.  On the forward half-plane
		$x=\operatorname{Re}(\bar\xi_1w)\geq0$, the latter alternative is
		impossible because $u^2\leq|w|^2+|P_1|^2\leq8$.  The cosine rule
		therefore gives
		\[
		2-x\leq e+5e^2\leq\frac54e,
		\qquad |w|\leq\frac25+\frac45x.
		\]
		Suppose $R_2\leq c$.  On the opposite half-plane the selected
		maximizer bounds $|w|$ by $c$; all points of radius below $c$ have the
		same bound.  Consequently the single globally convex test
		\[
		\psi(x)=\max\!\left\{c,\frac25+\frac45x\right\}^{n}
		\]
		majorizes $|w|^n$ on the whole convex hull. By \eqref{eq:boundary-candidate} and the fact that
		$F^*(e^{it})\in K$ for almost every $t$, we obtain
		\begin{equation}\label{eq:gaussian-single-end-cost}
			\begin{aligned}
				n|B_n(F)|
				&\leq \E_t |F^*(e^{it})|^n \\
				&\leq \E_t
				\psi\!\left(\operatorname{Re}(\bar\xi_1 F^*(e^{it}))\right) \\
				&\leq \E_t \psi(2\cos t).
			\end{aligned}
		\end{equation}
		The last inequality is the boundary projection comparison
		in the fixed direction $\xi_1$, applied to the convex function $\psi$.
		For $|\theta|\leq\pi$, the standard cosine parabola bound gives
		$\cos\theta\leq1-2\theta^2/\pi^2$.  Using
		$\max(0,1-v)^n\leq e^{-nv}$, the model integrand is bounded by
		\[
		\psi(2\cos\theta)
		\leq c^n+2^n\exp\!\left(-\frac{8n}{5\pi^2}\theta^2\right).
		\]
		Its normalized circle average is at most
		\[
		c^n+\frac{2^n}{2\pi}
		\sqrt{\frac{\pi}{8n/(5\pi^2)}}.
		\]
		Now $4n^2q^n\leq1$ for all $n\geq1501$: at $1501$ use
		$q^{40}<1/2$ and $4(1501)^2<2^{37}$, and propagate by
		$(1+1/n)^2q<1$.  Thus $q^n\leq1/(2000\sqrt n)$.
		Also $\pi<3.15$ gives
		$\sqrt{5\pi/(8n)}/2<141/(200\sqrt n)$.  The preceding cost,
		divided by $2^n$, is therefore less than
		$1411/(2000\sqrt n)$.  But
		$1411^2k<1000^2(2k-1)$ for $k\geq751$, so this is below
		$1/(2\sqrt k)\leq c_k=n\Cat_{k-1}/2^n$.
		This contradicts the candidate condition in
		\eqref{eq:gaussian-single-end-cost}.  The selected $P_2$ must have
		$R_2>c$.
	\end{proof}
	
	For the two-sided deficit argument, use the sharper
	$\delta=1/20$ and $p_0=39/20$.  If $|w|=2-e\geq p_0$
	and $|P_i|\geq|w|$, the two-point inequality and the
	parallelogram law give
	\[
	u(4-u)\leq16e,\qquad u=|w-P_i|.
	\]
	Since $e\leq1/20$, either $u\leq17e/4$ or
	$u\geq4-17e/4$.  In the forward half-plane of $P_1$
	the second alternative contradicts $u^2\leq8$.
	With $x=\operatorname{Re}(\bar\xi_1w)$, the cosine rule then gives
	\begin{equation}\label{eq:coarse-radius}
		2-x\leq e+5e^2.
	\end{equation}
	The opposite selected endpoint has $R_2>p_0$ by
	Lemma~\ref{lem:second-radius-195}.  For a high-radius point $w$
	in its assigned half-plane, the far alternatives relative to $P_1$
	give $|w+P_1|^2\leq17e$ and
	$|P_2+P_1|^2\leq17d_2\leq17e$.  Hence
	$|w-P_2|\leq2\sqrt{17e}<4-17/80$.
	The far alternative relative to $P_2$ is impossible, so
	\eqref{eq:coarse-radius} holds in direction $\xi_2$ as well.
	On each assigned half-plane, $R_i$ caps the radius of
	$\operatorname{conv}K$; the cutting line itself has radius below
	$1.9$ by the two-point inequality.  Thus the same convex test
	applies to both assigned sides.
	
	Define, for $x\leq2$,
	\[
	h(x)=\max\left\{p_0,\,
	2-\frac{\sqrt{1+20(2-x)}-1}{10}\right\},
	\qquad f(x)=h(x)^n,
	\]
	and extend $h(x)=x$ for $x\geq2$.  The inverse in this
	formula solves $s=e+5e^2$.  Consequently \eqref{eq:coarse-radius}
	and the radius caps give $|w|^n\leq f(\operatorname{Re}
	(\bar\xi_iw))$ on the assigned side.  Put
	\[
	V=\E f(2\cos\theta),\qquad
	\alpha=\arccos(31/32),\qquad
	\alpha_*=\arccos(39/40).
	\]
	The active part of the positive semicircle is $0\leq\theta\leq\alpha$;
	the rest is charged at most $p_0^n$.  The substitution
	\[
	e(v)=2-2\cos v,\qquad
	\Phi(v)=2\arcsin\!\left(\sin(v/2)\sqrt{1+5e(v)}\right),
	\quad 0\leq v\leq\alpha_*
	\]
	has $\Phi(0)=0$, $\Phi(\alpha_*)=\alpha$, and
	$h(2\cos\Phi(v))=2\cos v$.  Its derivative is
	\[
	\Phi'(v)=
	\frac{1+10e(v)}
	{\sqrt{(1+5e(v))(1-5e(v)^2/(4-e(v)))}}.
	\]
	The denominator is at least one: after clearing positive factors
	this is $4-2e-5e^2\geq0$ for $0\leq e\leq1/20$.
	Also $e(v)\leq(80/79)\sin^2v$.  With
	$C_n=\int_0^{\pi/2}\cos^n v\,dv$,
	integration by parts gives
	$\int_0^{\pi/2}\cos^n v\sin^2v\,dv=C_n/(n+2)$
	and $2^n C_n/\pi=I_n/2$.  Thus
	\begin{equation}\label{eq:coarse-model}
		V\leq p_0^n+\frac{I_n}{2}
		\left(1+\frac{800}{79(n+2)}\right)
		\leq p_0^n+\frac{I_n}{2}(1+11/n).
	\end{equation}
	The first integral is a regular angle-to-angle change of variables;
	$\Phi'(0)=1$, so no endpoint density singularity is introduced.
	
	Suppose a selected side has $d_i\geq1/n$.  Set
	$b=\sqrt{3/n}$, $x_0=2\cos b$, and
	$\mu=2\sin b/b\geq2-1/n$.  Replace $f$ to the right of
	$x_0$ by the affine line through $(x_0,f(x_0))$ and
	$(\mu,f(\mu))$.  The resulting convex test still controls
	the capped actual side; call its circle cost $T$.
	Writing $E_b g=b^{-1}\int_0^b g(\theta)\,d\theta$ and
	$X(\theta)=2\cos\theta$, affine averaging gives the exact identity
	\begin{equation}\label{eq:chord-arc-exact}
		V-T=\frac b\pi\bigl(E_b f(X)-f(E_bX)\bigr),
		\qquad E_bX=\mu .
	\end{equation}
	The line and $f$ may cross inside the arc; the identity uses their
	averages, not a pointwise saving.
	
	Here is a quantitative Jensen bound using only one-dimensional
	integrals.  On $[0,b]$, put $S=2-X$.  The cosine Taylor bounds give
	$(999/1000)\theta^2\leq S\leq\theta^2$.  Since
	$E_b\theta^2=b^2/3$ and $E_b\theta^4=b^4/5$,
	\[
	E_bX^2-(E_bX)^2=E_bS^2-(E_bS)^2
	\geq\frac{99}{100}\frac{4b^4}{45}.
	\]
	On $[x_0,2]$ the smooth branch of $h$ satisfies
	$h\geq2-3/n$, $h''\geq0$, and
	$(h')^2=(1+20(2-x))^{-1}$.  The logarithmic bound for
	$(1-3/(2n))^{n-2}$ and $e^{3/2}<9/2$ then give
	\[
	f''(x)\geq m_n:=2^{n-2}n(n-1)/5.
	\]
	Thus $f(x)-m_nx^2/2$ is convex on the arc range.  Ordinary
	integral Jensen and \eqref{eq:chord-arc-exact} yield
	\[
	V-T\geq\frac b\pi\frac{m_n}{2}
	\frac{99}{100}\frac{4b^4}{45}.
	\]
	The odd cosine-moment recurrence gives $I_n\leq2^n/\sqrt n$.
	After $b^2=3/n$, the last display exceeds
	\[
	I_n\,\frac{\sqrt3}{50\pi}\frac{99}{100}
	(1-1/n)>I_n/100
	\qquad(n\geq1501),
	\]
	where the final strict inequality follows from
	$\sqrt3>17/10$, $\pi<22/7$, and $1-1/n>99/100$.
	
	The platform estimate $2p_0^n\leq I_n/n$ follows from
	$(39/40)^{40}<1/2$ at $n=1501$ and the decreasing
	adjacent-term ratio.  The two unchanged side costs total at most
	$I_n(1+12/n)$ by \eqref{eq:coarse-model}.  If one side is bad,
	the chord saves more than $I_n/100$; the resulting cost is strictly
	below $(1-1/n)I_n$ because $13/n<1/100$.
	This contradicts the candidate lower moment in
	\eqref{eq:boundary-candidate}.  Both selected actual endpoints
	therefore satisfy
	\[
	0\leq nd_1<1,\qquad 0\leq nd_2<1.
	\]
	For the following positive-deficit estimates, a zero deficit is
	handled separately by slit rigidity following \eqref{eq:two-point}.
	For $d_i>0$, the localization gives $d_i<1/n\leq1/1501$.
	Apply the actual root strip \eqref{eq:whole-strip} at
	$r=1+\sqrt{d_i}/10$.  Write $L_i=\log(2/d_i)$.
	The principal width term is at most $(11/50)\sqrt{d_i}$.
	For the odd-tail kernel $S_r$, put
	$q=(r^2+1)/((r+1)/10)$.  Then
	$S_r-L_i/4=(\log q)/2-(\log2)/4-r^{-2}$.
	Here $r\leq1+1/380$, so $q/10-1<1/500$ and
	$r^{-2}>99/100$.  The inequality
	$\log(q/10)\leq q/10-1$, together with
	$\log5<13/8$ and $\log2<7/10$,
	shows $S_r\leq L_i/4$.
	Also $2/d_i>2048=2^{11}$ and
	$\log2>69/100$, whence $L_i>15/2$ and
	$\sqrt{L_i}>27/10$.  Since the odd-root energy factor
	$1-R_i^2/4$ is at most $d_i$, the squared strip width satisfies
	$W_i\leq d_iL_i(11/135+1/2)^2<0.34d_iL_i$.
	Thus
	\begin{equation}\label{eq:actual-width-034}
		W_i\leq0.34d_iL_i,
		\qquad y_i:=nW_i\leq0.34t_i\log(2n/t_i).
	\end{equation}
	The inequality holds on the whole ordinary convex hull, not just at
	its two selected endpoints.
	For the sharper cost estimate, put $u=1501d_i\leq1$.  The elementary
	bound $-\log u\leq2(u^{-1/2}-1)$ and
	$2<\log3002<25/3$ give
	\[
	\sqrt{1501d_i}\,\log(2/d_i)<25/3.
	\]
	Consequently the unified upper width
	$\widehat W_i=0.34d_i\log(2/d_i)$ satisfies
	\[
	0<W_i<\widehat W_i<1/500,\qquad
	n\widehat W_i^2<3/500.
	\]
	Here the last inequality uses $nd_i<1$; it holds uniformly as
	$n$ increases.  These two scalar bounds and their transfer to the
	actual squared strip width follow from the same inequalities.
	
	The finer model below uses the lower plateau $1.9$.  To assign every
	point above that plateau, repeat the two-point argument with
	$0\leq e=2-|w|\leq0.1$ and
	$C=16/(2+\sqrt{4-16/10})<4.51$.
	The alternatives are $|w-P_i|\leq Ce$ and
	$|w-P_i|\geq4-Ce$.  The forward half of $P_1$ permits only the
	near alternative.  In the opposite half, $w$ and the selected $P_2$
	are both far from $P_1$, so the two-point inequality gives
	$|w+P_1|\leq2\sqrt{Ce}$ and
	$|P_2+P_1|\leq2\sqrt{Cd_2}$.
	As $d_2\leq e$, their difference has length at most
	$4\sqrt{Ce}<4-Ce$; thus it is near $P_2$.
	The cosine rule now shows that each assigned point above $1.9$
	has positive projection in its own endpoint direction.
	For a positive-deficit endpoint, $W_i<1/100$ follows already from
	\eqref{eq:actual-width-034} and $d_i<1/1501$.
	On a nonnegative projection $x\leq2$, the squared strip implies
	$|w|\leq\max\{1.9,2-(1-1.07W_i)(2-x)+1.07W_i^2\}$:
	for $x\leq949/500$ the strip square is below $1.9^2$;
	above $949/500$, expanding the displayed affine square dominates
	the strip square.  The latter comparison only uses
	$(2.14)x>4$ and $W_i\geq0$.
	
	For the final model bound put $x=\operatorname{Re}(\bar\xi_iw)$,
	$s=2-x$, and use
	\[
	h_i(2-s)=\max\{1.9,
	2-(1-1.07W_i)s+1.07W_i^2\}.
	\]
	On its active arc $s<0.102$.  The squared strip
	\eqref{eq:whole-strip} makes $h_i$ a radius majorant there;
	outside, the plateau applies.  Set
	$U=1+1.07W_i^2/2$, $v=1-1.07W_i$, and $\gamma=v/U$.
	If $s=2-2\cos\theta$ is the original angular deficit, choose
	$\phi$ by $2-2\cos\phi=\gamma s$.  The active branch is exactly
	$2U\cos\phi$, and the angular Jacobian is
	\[
	\frac{d\theta}{d\phi}
	=\gamma^{-1/2}
	\sqrt{\frac{4-\gamma s}{4-s}}.
	\]
	Since $1-\gamma\leq1.07107W_i$ and $s\leq0.102$, the square-root
	factor is at most $1+0.0142W_i$.  The remaining active model factor is
	\[
	(1+0.0142W_i)(1+1.07W_i^2/2)^{n+1/2}
	(1-1.07W_i)^{-1/2}
	\leq1+(0.57+0.55y_i)W_i.
	\]
	For this last inequality, take logarithms and use
	$W_i<1/500$ and $nW_i^2<3/500$; the logarithm is at most
	$0.555W_i+0.536nW_i^2<1/200$, and
	$e^x\leq1+(200/199)x$ on $0\leq x\leq1/200$.
	With the platform charged once, the resulting one-end cost is
	\[
	V_i\leq(1.9)^n+\frac{I_n}{2}
	\bigl(1+0.57W_i+0.55nW_i^2\bigr).
	\]
	This bound has also been applied to both selected actual endpoints
	in the candidate inverse-coefficient inequality.
	For the refined affine $h_i$, use the same chord with
	$b=\sqrt{3d_i}$ and $\mu=2\sin b/b\geq2-d_i$;
	the resulting convex extension majorizes the assigned side because
	$h_i(\mu)\geq\mu\geq2-d_i$.  Its top cap saves at least
	$0.067I_nt_i^{5/2}e^{-0.61t_i}$ on side $i$.
	To see the size of this saving, use the normalized average
	$E_b$ on $[0,b]$ as in \eqref{eq:chord-arc-exact}. Set
	\[
	S=2-2\cos\theta,\qquad P(s)=(2U-vs)^n,\qquad
	A=n(n-1)2^{n-2}v^2,\quad B=n/2,\quad s_0=2b^2/5 .
	\]
	The entire short arc lies on the affine branch: $S\leq b^2<3/n$
	and $2U-vS\geq2-S>1.9$. The chord identity therefore gives
	$V_i-T_i=(b/\pi)\{E_bP(S)-P(E_bS)\}$, where $T_i$ is the
	cost after replacing the cap by its chord.
	
	For $0\leq s\leq b^2$, the inequality
	$\log(1-s/2)\geq-s/(2-s)$ and $ns<4$ give
	$(1-s/2)^{n-2}\geq e^{-ns/2}$. Thus $P''(s)\geq A e^{-Bs}$.
	The tangent to this exponential at $s_0$ shows that
	$P-Ae^{-Bs_0}Q$ is convex, where
	\[
	Q(s)=\tfrac12(1+Bs_0)s^2-\tfrac16Bs^3 .
	\]
	To apply Jensen, put $X=S/b^2$, $u=\theta/b$ and $c=999/1000$.
	The cosine Taylor bounds give $cu^2\leq X\leq u^2$, since
	$b^2<3/1501$. As $E_bu^j=1/(j+1)$, we have
	\begin{align*}
		\operatorname{Var}_b(S)/b^4
		&\geq c^2/5-1/9> (199/200)(4/45),\\
		\{E_b S^3-(E_b S)^3\}/b^6
		&\leq 1/7-c^3/27<(6/5)(199/200)(4/45).
	\end{align*}
	Since $3s_0=6b^2/5$, the cubic difference is at most
	$3s_0\operatorname{Var}_b(S)$. In particular
	\[
	E_bQ(S)-Q(E_b S)
	=\tfrac12(1+Bs_0)\operatorname{Var}_b(S)
	-\tfrac16B\{E_bS^3-(E_b S)^3\}
	\geq\tfrac12\operatorname{Var}_b(S).
	\]
	Jensen now gives $V_i-T_i\geq (b/\pi)Ae^{-Bs_0}
	\operatorname{Var}_b(S)/2$. Substituting $b^2=3t_i/n$ and using
	$I_n<(4/5)2^n/\sqrt n$ (from the product of consecutive cosine
	moments) gives
	\[
	\frac{V_i-T_i}{I_n t_i^{5/2}e^{-3t_i/5}}
	>\frac{\sqrt3}{8\pi}\frac{199}{200}(1-1/n)v^2
	>\frac{67}{1000}.
	\]
	For the last step one may substitute $\sqrt3>173/100$,
	$\pi<22/7$, $n\geq1501$, and $v\geq1-1.07/500$.
	This proves the stated saving since $e^{-3t_i/5}\geq e^{-0.61t_i}$.
	Summing the two side costs and comparing with
	\eqref{eq:boundary-candidate} yields
	\begin{equation}\label{eq:actual-joint-J}
		J_n(t_1)+J_n(t_2)
		\leq0.502+\frac{2n(1.9)^n}{I_n}<0.52,
	\end{equation}
	where the definition of $J_n$ uses the upper $y_i$ in
	\eqref{eq:actual-width-034}.  The preceding bound $J_n(t)>-9/160$ then gives $J_n(t_i)<0.58$ separately for both sides.
	
	\begin{lemma}[Cutoff from the scalar input]\label{lem:cutoff}
		If $n\geq1501$, $0<t<1$, and $J_n(t)<0.58$, then
		\[
		t<\frac3{20},\qquad d<\frac1{10000},\qquad
		tL<\frac32,\qquad \Delta L<\frac32.
		\]
	\end{lemma}
	\begin{proof}
		Write $t_0=3/20$, $n_0=1501$, $z_*=0.51$, and
		$L_0=\log(2n_0/t_0)<10$.  A direct rational comparison gives
		\[
		0.067n_0t_0^{5/2}e^{-0.61t_0}
		-0.285z_*-0.275z_*^2>0.58.
		\]
		For example, use $\sqrt{t_0}>12/31$ and
		$e^{-0.61t_0}\geq(1-183/4000)^2$.
		For $t_0\leq t\leq1$, the factor
		$\sqrt t\,e^{-0.61t}$ is at least its value at $t_0$.
		On $[t_0,3/5]$ this follows from
		$\log(t/t_0)\geq1-t_0/t\geq(61/50)(t-t_0)$;
		on $[3/5,1]$ the estimates
		$\sqrt t>3/4$ and $e^{-0.61t}>1/2$ suffice,
		because $\sqrt{t_0}<0.39$ and $e^{-0.61t_0}<0.94$.
		Also $L/n\leq L_0/n_0$ and $L^2/n\leq L_0^2/n_0$:
		first decrease $t$ to $t_0$, then use the decreasing functions
		$\log x/x$ and $\log x/\sqrt x$ for
		$x=2n/t_0>e^2$.  Divide $J_n(t)$ by $nt^2$ and
		compare all three terms with the base point.  The base inequality
		then gives $J_n(t)>0.58$ throughout $t\geq t_0$.
		
		For $0<t\leq t_0$, apply the cubic envelope from the preceding
		paragraph with $t_0$ in place of $1$.  The base checks
		$\sqrt{t_0}<0.39$ and
		$L_0<9.91$ give
		\[
		100(0.34)^3\sqrt{t_0}L_0^3<1501,
		\qquad 100z^3\leq nt^2\sqrt t.
		\]
		One may check $L_0<9.91$ from
		$e^{9.91}=e^{10}e^{-0.09}>
		(2.718)^{10}(0.91)>2n_0/t_0$.
		Since $e^{-0.61t}\geq(1-183/4000)^2$, the positive
		term of $J_n$ exceeds $6.1z^3$.  A polynomial expansion
		about $z=0.51$ shows
		\[
		6.1z^3-0.285z-0.275z^2>0.58
		\quad(z\geq0.51).
		\]
		Thus the candidate condition forces $z<0.51$.
		As $z=0.34tL$, this gives $tL<1.5$; also
		$t<t_0$ gives $d=t/n<1/10000$, and $m<n$
		gives $\Delta L<tL<1.5$.
	\end{proof}
	
	For the entire hull use the root strip at $r=1+(2/25)\sqrt d$.
	Here $d\leq1/5000$ and $L>9$. Set $u=\sqrt d<1/70$ and
	$q=(r^2+1)/((2/25)(r+1))$. The formulas in the root estimate give
	\[
	\frac{p_r}{u}
	=\frac{4/25+(1/2+4/625)u}{1+2u/25}<\frac{17}{100},
	\qquad
	S_r-\frac L4=\frac12\log q-\frac14\log2-r^{-2}.
	\]
	For the second formula, $q/(25/2)-1=r(r-1)/(r+1)<1/875$
	and $r^{-2}\geq1-2(r-1)>1-2/875$. Using $\log(1+x)\leq x$,
	$\log5<161/100$ and $\log2>69/100$ shows
	\[
	S_r-\frac L4<
	\frac{161}{100}-\frac34\frac{69}{100}+\frac1{1750}
	-\left(1-\frac2{875}\right)<\frac1{10}.
	\]
	The logarithm bounds follow from the Taylor sum through degree
	eight for $e^{161/100}$, which exceeds $5$, and
	$\log2>2(1/3+1/81)$. By Lemma~\ref{lem:cutoff}, $d<1/10000$ and $2/d>2^{14}$,
	so $L>9$. Since $E\leq d$ and $1/4+1/90<(13/25)^2$, we obtain
	\begin{equation}\label{eq:width-input}
		\frac{W}{dL}\leq
		\left(\frac{17}{100\sqrt L}+\sqrt{\frac14+\frac1{10L}}\right)^2
		<\left(\frac{173}{300}\right)^2<\frac13 .
	\end{equation}
	Define, only for $d>0$,
	\begin{equation}\label{eq:new-a}
		y=\Delta L/3,\qquad a=\min\{3/y,m/5\}.
	\end{equation}
	Then $mW<y<1/2$, $a>6$, and $ay\leq3$.
	
	We also need positive projection on each assigned outer half.
	In first-axis coordinates write $p+iq=\bar\xi_1w$.
	The first whole-hull strip has $W_1<1/3000$, because $mW_1<1/2$.
	Together with $|w|\leq2$ and $R_1\leq2$ it gives $q^2<1/100$
	on the entire hull, including its negative half.
	If $|w|\geq46/25$, then $|p|>183/100$.
	The selected second endpoint has $R_2>1999/1000$, negative first
	projection, and the same transverse bound, so its first projection
	is less than $-199/100$.
	Consequently, on the negative assigned half,
	\[
	\operatorname{Re}(\bar\xi_2w)>
	\frac{(199/100)(183/100)-1/100}{2}>\frac{181}{100}.
	\]
	On the positive half the first projection already exceeds $183/100$.
	A change of unit direction by at most $1/200$ alters projection by
	at most $1/100$.  The nearest phase ray has distance at most
	$\pi/(2k)<1/200$, so both assigned outer projections exceed $9/5$.
	The whole-hull cap under tilting is proved below; it uses the same
	selected pair and applies on every source circle.
	
	\subsection{The saving from a capped projection}
	Write $T_m=2^m m^{-3/2}$. We compare the projected filling on each
	circle with the Poisson extension of the flattened cosine. Only the
	outer annulus will be needed to obtain a strict gain after integration.
	
	\begin{lemma}[Short actual disk gain]\label{lem:actual-disk-gain}
		Let $m\geq1500$, $d>0$, and $D=md\leq3/20$. Suppose $|\xi|=1$ and
		$\sup_{w\in K}\operatorname{Re}(\bar\xi w)\leq 2-d$. Then
		\[
		\frac{m\sqrt m}{2^m}\int_{\D}
		\left\{(2\operatorname{Re}z)_+^m
		-\bigl(\operatorname{Re}(\bar\xi H(z))\bigr)_+^m\right\}d\sigma
		>\frac{D^3}{30}.
		\]
		
	\end{lemma}
	\begin{proof}
		The boundary projection $u(t)=\operatorname{Re}(\bar\xi F^*(e^{it}))$
		satisfies $u\leq2-d$ almost everywhere since the trace lies in $K$.
		Write $2\sin h/h=2-d$, $0<h<\pi$, and replace $2\cos t$ on $[-h,h]$
		by its mean $2-d$ to obtain $u_d$. For $0<r<1$, write
		$Y_r(t)=\operatorname{Re}(\bar\xi H(re^{it}))=(P_r*u)(t)$.
		The boundary comparison for $u_d$ in Section~\ref{sec:bridge} and
		\eqref{eq:poisson-convex-order} give
		$Y_r\preceq_{\rm cx}P_r*u_d$ for normalized angular measure.
		Put $M(r)=(P_r*u_d)(0)$. Each superlevel set of $u_d+2$ is a centered
		arc, over which the even, decreasing Poisson kernel has its largest
		integral when centered at zero. Integration over levels gives
		$(P_r*u_d)(t)\leq M(r)$ for every $t$.
		Now test the convex order with $x\mapsto(x-M(r))_+$:
		\[
		0\leq \E_t (Y_r(t)-M(r))_+
		\leq \E_t\bigl((P_r*u_d)(t)-M(r)\bigr)_+=0 .
		\]
		Since $Y_r$ is continuous, $Y_r(t)\leq M(r)$ at every angle.
		The cap angle obeys $\sqrt{3d}\leq h\leq(7/4)\sqrt d$.
		For the first bound use $\sin h\geq h-h^3/6$. For the second,
		$2-2\sin v/v$ increases on $(0,\pi)$, and at $v=(7/4)\sqrt d$
		the sine Taylor bound gives $2-2\sin v/v\geq v^2/3-v^4/60>d$,
		using $d\leq1/10000$.
		Let $P_r(\theta)=(1-r^2)/(1-2r\cos\theta+r^2)$ and
		$A_r=\pi^{-1}\int_0^hP_r(\theta)\,d\theta$.
		On $[0,h]$ the cap difference times the kernel is
		$(d+(1-r)^2/r)P_r(\theta)-(1-r^2)/r$. Integration yields
		\[
		2r-M(r)=\left(d+\frac{(1-r)^2}{r}\right)A_r
		-\frac{(1-r^2)h}{\pi r}.
		\]
		The Poisson kernel has complementary mass
		$1-A_r\leq(1-r^2)/(\pi r h)$: integrate its elementary
		$\sin^{-2}(\theta/2)$ upper bound over $[h,\pi]$ and use
		$\cot(h/2)\leq2/h$. Dropping the nonnegative term
		$(1-r)^2A_r/r$ bounds the defect below by
		$d-(1-r^2)(d/h+h)/(\pi r)$. For $1-\sqrt d/3\leq r<1$
		we have $r\geq99/100$ and
		$((1+r)/(\pi r))(d/h+h)\leq(3/2)\sqrt d$, using
		$1/\sqrt3<7/12$, $h\leq(7/4)\sqrt d$, and $\pi>157/50$.
		Thus
		\begin{equation}\label{eq:short-peak-defect}
			2r-M(r)\geq d-\tfrac32(1-r)\sqrt d>0.
		\end{equation}
		
		Let $f(x)=x_+^m$ and let $\psi_M$ agree with $f$ below $M$ and
		continue by its tangent above $M$. Taking $M=M(r)$ gives
		$f(Y_r)=\psi_M(Y_r)$ pointwise. Since $\psi_M$ is convex,
		\eqref{eq:ordinary-convex} yields
		$\E_t f(Y_r(t))\leq\E_t\psi_M(2r\cos t)$ on each source circle.
		Put $v=d-\tfrac32(1-r)\sqrt d$ and $q=2r-v$.  On
		$|\theta|\leq\sqrt v$ one has $2r\cos\theta\geq q\geq M(r)$.
		The binomial remainder above a tangent therefore gives
		\[
		f(2r\cos\theta)-\psi_M(2r\cos\theta)
		\geq\frac{m(m-1)}2q^{m-2}(v-\theta^2)^2.
		\]
		As $\int_0^{\sqrt v}(v-\theta^2)^2d\theta
		=(8/15)v^{5/2}$, the full normalized circle saves at least
		$4m(m-1)q^{m-2}v^{5/2}/(15\pi)$.
		
		Only the outer annulus $a\leq r<1$, $a=1-\sqrt d/16$, is needed.
		There $v\geq(9/10)d$, $r\geq99/100$, and
		$(1-v/(2r))^{m-2}>9/10$.  Hence
		$v^{5/2}\geq(3/4)d^{5/2}$ and the circle saving is at least
		\[
		\frac9{50\pi}m(m-1)(2r)^{m-2}d^{5/2}.
		\]
		The radial factor is $2r\,dr$, so its integral contributes exactly
		$2^{m-1}(1-a^m)/m$.  After multiplication by $m\sqrt m/2^m$,
		the saving is at least
		\[
		\frac9{100\pi}\left(1-\frac1m\right)
		D^2\sqrt D\,(1-a^m).
		\]
		Writing $x=\sqrt D\leq2/5$ gives $a^m\leq
		e^{-\sqrt m\,x/16}\leq e^{-12x/5}$. Concavity bounds
		$1-e^{-12x/5}$ below by the chord from $0$ to $2/5$.
		Its slope exceeds $4/3$ since
		$e^{24/25}>1+24/25+(24/25)^2/2>15/7$.
		Using $\pi<22/7$ and $1-1/m\geq99/100$, the last display is
		at least $(189/5000)D^3>D^3/30$.  On every other source circle
		the ordinary positive-power convex comparison makes the radial
		increment nonnegative.  Integrability and exhaustion by interior
		disks pass the strict inequality to the open unit disk.
	\end{proof}
	
	\paragraph{Tilting the second cap.}
	Write $P_i=R_i\xi_i$ for the same selected endpoints.  Since $P_1$
	maximizes the modulus on the entire hull, every unit direction has
	projection at most $R_1$.  For the second cap, split the hull by
	$\operatorname{Re}(\bar\xi_1w)=0$.  On the negative half its
	modulus is at most $R_2$.  On the positive half, the fact that
	$\operatorname{Re}(\bar\xi_1\xi_2)\leq0$ gives
	$\operatorname{Re}(\bar\xi_2w)\leq
	|\operatorname{Im}(\bar\xi_1w)|<1/10$ by the first strip.
	For $|v-\xi_2|\leq1/4$ and $|w|\leq2$, the perturbation of this
	projection is at most $1/2$.  Hence it remains below $R_2>1$.
	Thus the whole-hull cap, and the actual saving above, hold along
	every such tilted unit direction.  All width and maximality data
	refer to the original selected pair.
	
	For any unit coefficient phase $\beta$, the ray can be chosen
	explicitly:
	\[
	v=\xi\exp\!\left(-\frac{i}{2k}
	\operatorname{Arg}((\beta\xi^k)^2)\right).
	\]
	Then $(\beta v^k)^2=1$, so the signed angular loss vanishes on
	this ray, and $|v-\xi|\leq\pi/(2k)<1/4$.  Squaring the phase
	chooses either sign of a sine zero and includes the tie case.
	
	\begin{lemma}[Retaining a power saving]\label{lem:two-power-retention}
		Let $m\geq1500$, $\ell=m-8$, $c\geq0$, and
		$\kappa=\ell(\ell-1)/(m(m-1))$.  Put
		$g(u)=\max\{u_+^m,c^8u_+^\ell\}$.
		The function $g(u)-\kappa u_+^m$ is convex on the whole real line.
		For every actual filling and unit direction $v$, therefore,
		\[
		\int_{\D}\!\left[g(\operatorname{Re}z)
		-g\!\left(\tfrac12\operatorname{Re}(\bar vH)\right)\right]d\sigma
		\geq\kappa\int_{\D}\!\left[(\operatorname{Re}z)_+^m
		-\left(\tfrac12\operatorname{Re}(\bar vH)\right)_+^m\right]d\sigma.
		\]
		Under the cap of Lemma~\ref{lem:actual-disk-gain}, the left side,
		multiplied by $m\sqrt m$, is strictly greater than $D^3/32$.
	\end{lemma}
	\begin{proof}
		Below zero the residual is zero.  On $[0,c]$ its second
		derivative is
		$\ell(\ell-1)u^{\ell-2}(c^8-u^8)\geq0$; above $c$ the residual
		is $(1-\kappa)u^m$.  The derivative jumps upward by
		$8c^{m-1}$ at $c$, and the derivatives at zero agree.
		This proves global convexity, including $c=0$.
		Apply the actual disk convex comparison to this residual and
		rearrange.  All tests are continuous on the bounded projection
		range, so their integrals exist.  The final assertion follows
		from $\kappa\geq49/50$ and
		$(49/50)/30>1/32$, with the same direction in both integrals.
	\end{proof}
	
	\paragraph{The extremely small deficit regime.}
	The following comparison resolves the joint two-side cost once the
	actual endpoint localization has supplied the width and cap hypotheses.
	Let $d=2-\max_K|w|>0$, put $m=2k-2$ and $x=\log(100m)$, and assume
	\begin{equation}\label{eq:tiny-cutoff}
		d\leq d_0:=\frac{3}{100m^2x}.
	\end{equation}
	Since $d$ is the smallest endpoint deficit, every projection of
	the boundary trace has the global cap $2-d$.  The strip formula
	\eqref{eq:whole-strip} together with the endpoint width
	$W\leq dL/3$ gives, on the whole hull,
	\[
	|\operatorname{Im}(\bar\xi w)|<2.5\sqrt{dL},
	\qquad L=\log(2/d),
	\]
	where $\xi$ is the direction of a maximal-radius point.  Indeed
	$|\operatorname{Im}(\bar\xi w)|^2\leq4W(W+4)$ and
	$dL$ is tiny in \eqref{eq:tiny-cutoff}.
	Choose the nearest pair of antipodal phase rays $v,-v$ to
	$\xi,-\xi$, and set $b'=6k\sqrt{dL}$.  Assign $w$ by the sign
	of its true-axis projection.  On $|w|>b'$ that positive projection
	exceeds $|w|/2$, since the transverse square is at most
	$(25/4)dL$.  The argument bound $|\arg u|\leq|\operatorname{Im}u|/
	\operatorname{Re}u$ therefore gives a scaled true-axis angle at
	most $5/6$.  The nearest phase ray adds at most $\pi/2$, giving
	a scaled phase angle below $2.8$.
	For $0\leq t\leq2.8$, concavity and the cubic sine bound at
	$\pi-2.8$ give $\sin t\geq t/10$.  Together with Bernoulli's
	inequality and $m/2=k-1\leq 3k^2/1400$ (since $k\geq751$), this gives
	\[
	1-\frac3{14}\sin^2(k\theta)\leq\cos^m\theta
	\qquad(|k\theta|\leq2.8).
	\]
	Thus the angular penalty pays the projected-power loss outside
	the central region $|w|\leq b'$.
	Inside, a residual at most $(b'/2)^m$ remains in the dimensionless
	signed integral.
	
	Each of the two projected rays has cap $2-d$, so the full convex
	capped-Poisson comparison applies to its positive $m$th-power test.
	In the dimensionless normalization of \eqref{eq:signed-criterion}, each
	ray saves more than $\Delta^3/(30m^{3/2})$ from its reference
	moment $M_m$, by Lemma~\ref{lem:actual-disk-gain}.  The two rays therefore save more than
	$\Delta^3/(15m^{3/2})$.
	It remains to pay the variable central residual.  This can be
	done without differentiating a ratio.  Put $L_0=\log(2/d_0)$.
	The bounds $1\leq x\leq m/20$ give $1\leq L_0\leq3x$.
	For $d\leq d_0$, the logarithmic tangent inequality gives
	\[
	dL=dL_0+d\log(d_0/d)\leq dL_0+d_0-d\leq d_0L_0.
	\]
	Consequently $m^2dL\leq9/100$, and $m\geq9k/5$ implies
	$9k^2dL\leq1/4$.  Also $dL^2\leq1$: the elementary Taylor
	bound $e^L\geq L^4/24\geq2L^2$ for $L\geq9$, together with
	$e^L=2/d$, proves this directly.
	Writing $q=k-1$, retain six powers and bound all the others by
	$1/4$:
	\[
	(b'/2)^m=(9k^2dL)^q
	\leq4^{-(k-7)}(9k^2)^6d^3(dL^2)^3
	\leq4^{-(k-7)}9^6k^{12}d^3.
	\]
	The sequence $15\cdot9^6 k^{12}4^{-(k-7)}$ is below one for
	$k\geq751$.  At the first order use $751<2^{10}$ and
	$15\cdot9^6<2^{24}$; successive ratios are at most
	$(752/751)^{12}/4<1$.  Since $m\sqrt m\leq m^3$, we obtain
	\begin{equation}\label{eq:tiny-ratio}
		15m\sqrt m\,(b'/2)^m<(md)^3.
	\end{equation}
	The signed cost is therefore strictly below $2M_m$ in this
	regime.  This is a joint estimate: no artificial lower cutoff on
	the tiny endpoint deficit is needed.  Its only actual-map inputs are
	the width and cap statements specified above.
	
	In the aligned phase-ray branch the desired final comparison is a
	retained saving greater than $\Delta^3/32$ against model and central
	errors totaling less than $13\Delta^3/500<\Delta^3/32$, in units of $T_m$.
	Here is the pointwise check with the shorter $a$, conditional on the
	actual endpoint and tilt estimates.  Put
	$s=m(1-x/2)$, $u=x/2$, and
	$h=|\operatorname{Im}(\bar v w)|/x$ on the outer region $x\geq1.8$.
	For a phase ray within scaled angle $\gamma\leq\pi/3$ of the true
	endpoint direction, write
	$\theta=|\arg(v/\xi_i)|\leq\pi/(3k)$,
	$q_\xi=\operatorname{Im}(\bar\xi_iw)$ and
	$Q=m|q_\xi|$.  Rotation gives
	$R_i-\operatorname{Re}(\bar\xi_iw)
	\leq2s/m+|q_\xi|\tan\theta$.
	With $mW\leq y<1/2$, the squared strip implies
	\[
	Q^2\leq4y(y+2s+Q\tan\theta),
	\qquad 0\leq s\leq m/10,
	\qquad m\tan\theta\leq21/10.
	\]
	Here $\cos\theta\geq999/1000$, so $x\cos\theta\geq7/4$.
	Since $y\leq1/2$, $m\geq1500$, and
	$h\leq\tan\theta+|q_\xi|/(x\cos\theta)$, we obtain the two
	simple rational bounds
	\[
	Q^2\leq1+8ys+\frac7{2500}Q,
	\qquad mh\leq\frac{21}{10}+\frac47Q.
	\]
	If $s\leq7a/5$, then $ys\leq21/5$ and the first inequality
	forces $Q<6$.  Thus $mh<387/70$ and
	$k|\arg(w/v)|\leq kh<14/5$.
	To pay for the projection, it suffices to use
	$\sin t\geq t/10$ on $[0,14/5]$.  Indeed, the cubic sine bound
	at $\pi-14/5$, with $3.14<\pi<3.15$, gives
	$\sin(14/5)\geq3/10$; concavity supplies the whole interval.
	Bernoulli's inequality and $1-\cos\varphi\leq\varphi^2/2$
	now give
	\[
	1-\cos^m\varphi\leq\frac m2\varphi^2
	\leq\frac3{14}\frac{k^2\varphi^2}{100}
	\leq\frac3{14}\sin^2(k\varphi),
	\qquad |k\varphi|\leq14/5.
	\]
	Consequently the signed angular loss pays for the first test $u^m$.
	
	If $s\geq7a/5$, use $Q\leq1+Q^2/4$ to absorb the small linear
	term, obtaining $Q^2\leq(51/50)(1+8ys)$.  Completing the square
	with $(Q-49/10)^2\geq0$ then gives
	\[
	\frac{m^2h^2}{2}\leq\frac{11}{2}+\frac73ys
	\leq\frac{11}{2}+\frac76s<8(s-a).
	\]
	The final strict inequality uses only $a\geq6$ and $s\geq7a/5$.
	Taking logarithms shows that the
	second test $(1-a/m)^8u^{m-8}$ dominates the entire unsigned
	cost $(|w|/2)^m$; it uses no phase saving.
	
	The junction of the two tests is at $s=a$.  Their convexity and
	retention of the actual saving are supplied by
	Lemma~\ref{lem:two-power-retention}.
	
	\begin{lemma}[One disk budget for two aligned ends]\label{lem:aligned-disk-budget}
		Suppose both selected rays belong to the aligned phase case and
		choose their parameters by \eqref{eq:new-a}.  Write
		$g_i(u)=\max\{u_+^m,(1-a_i/m)^8u_+^{m-8}\}$ and
		\[
		\mathcal R_i=\int_{\D}g_i(\operatorname{Re}z)\,d\sigma,
		\qquad
		\mathcal A_i=\int_{\D}g_i\!\left(\tfrac12
		\operatorname{Re}(\bar v_iH)\right)d\sigma.
		\]
		Then, with $N=m\sqrt m$,
		\[
		N\left(\int_{\D}S(H)\,d\sigma-2M_m\right)
		<N\left[(\mathcal R_1-M_m)+(\mathcal R_2-M_m)+(23/25)^m\right]
		-\frac{\Delta_1^3+\Delta_2^3}{32}.
		\]
	\end{lemma}
	\begin{proof}
		On an assigned outer point the signed density is bounded by its
		own test.  Both tests are nonnegative, so extending each test to
		the entire source disk preserves an upper bound.  The central
		region adds at most $(23/25)^m$, counted once.  Thus
		$\int S(H)\leq\mathcal A_1+\mathcal A_2+(23/25)^m$.
		The actual retained gains give
		$N(\mathcal R_i-\mathcal A_i)>\Delta_i^3/32$ for the same
		directions and parameters.  Adding and rearranging proves the
		claim.  All integrals exist by the bounded projection range and
		continuity of the tests; the signed density is integrable by the
		actual energy construction.  No boundary trace identification or
		assumption that a dividing level has measure zero is used here.
	\end{proof}
	
	For the other phase branch $\gamma\geq\pi/3$, use the true endpoint
	axis and $f_0(u)=\lambda(u_+)^{m-8}$, where
	$\lambda=(1-1/m)M_m/M_{m-8}$.
	The true axis keeps the geometric estimate simple.  Write $h$ for
	the signed angle of $w$ from that axis and $s=m(1-x/2)$, where
	$x=\operatorname{Re}(\bar\xi_iw)\geq9/5$.  The actual strip gives
	\[
	\tfrac12m^2\tan^2 h\leq\frac16+\frac58s,
	\qquad s\leq2\ \Longrightarrow\ |kh|\leq\frac67.
	\]
	Indeed $R_i-x\leq2s/m$, so the squared strip and $mW\leq y<1/2$
	give $m^2|\operatorname{Im}(\bar\xi_iw)|^2\leq 1+4s$.
	Since $\tan h=\operatorname{Im}(\bar\xi_iw)/x$ and $x\geq9/5$,
	the first bound follows from $(1+4s)/(2(9/5)^2)\leq1/6+5s/8$.
	For the second, use $|h|\leq|\tan h|$ and $k/m\leq751/1500$;
	at $s=2$ the resulting bound on $(kh)^2$ is below $(6/7)^2$.
	The moment recurrence $(q+4)M_{q+2}=(q+1)M_q$ gives the exact ratio
	\[
	\lambda=\frac{m-1}{m}\prod_{j=0}^3
	\frac{m-2j-1}{m-2j+2},\qquad
	\lambda M_{m-8}=(1-1/m)M_m.
	\]
	Here the recurrence is used only at positive even exponents.
	Applying $\log(1-v)\geq-v/(1-v)$ to the five factors gives
	\[
	-m\log\lambda\leq\frac{m}{m-1}+3m\sum_{j=0}^3\frac1{m-2j-1}
	\leq\frac{m}{m-1}+\frac{12m}{m-7}<14,
	\]
	using the decreasing ratios at $m=1500$.  If $\gamma$ is the
	signed phase displacement, then $\pi/3\leq|\gamma|\leq\pi/2$.
	Indeed the logarithm of the signed density divided by
	$\lambda(x/2)^{m-8}$, multiplied by $m$, is
	\[
	m\log\!\left(1-\tfrac3{14}\sin^2(\gamma-kh)\right)
	+8m\log(x/2)+\tfrac12m^2\log(1+\tan^2 h)-m\log\lambda .
	\]
	Use $x/2=1-s/m$, $|w|^2/x^2=1+\tan^2h$, and $\log(1+v)\leq v$
	for $v>-1$. The resulting upper bound is
	\begin{equation}\label{eq:true-axis-corrected}
		14+\frac16+\frac58s-8s-
		\frac3{14}m\sin^2(\gamma-kh).
	\end{equation}
	For $s\geq2$ the first four terms already have negative sum.
	For $0\leq s\leq1$, the strip estimate above gives
	\[
	|kh|^2\leq (k/m)^2(1/3+5s/4)
	\leq (751/1500)^2(19/12)<4/9 .
	\]
	Thus $|\gamma-kh|$ lies between $\pi/3-2/3$ and $\pi/2+2/3$.
	Its sine is greater than $1/4$, since the distance of this interval
	from both $0$ and $\pi$ exceeds $1/3$ and $\sin(1/3)>1/4$.
	The penalty in \eqref{eq:true-axis-corrected} is then larger than
	$3m/(14\cdot16)>15$, whereas the other terms sum to at most $14+1/6$.
	For $1\leq s\leq2$ use $|kh|\leq6/7$.
	The minimum of the sine on $[\pi/3-6/7,\,\pi/2+6/7]$ is its
	value at $\pi/3-6/7$, which exceeds $3/16$ by the cubic Taylor
	bound. Consequently $\sin^2(\gamma-kh)>1/30$.
	Here the penalty is greater than $3m/(14\cdot30)>10$,
	while the other terms sum to at most $14+1/6+5/8-8<7$.
	This proves the strict true-axis comparison.  Its actual disk
	integral is at most $(1-1/m)M_m$ by convex projection comparison.
	
	Thus each assigned outer half has its own test: the two-power
	test in the aligned case and the true-axis lower power in the
	other case.  Both are nonnegative.  Extending both to the entire
	source disk and adding the central allowance once gives
	$\int S(H)\leq\mathcal A_1+\mathcal A_2+(23/25)^m$ for every
	phase assignment.  This statement and the exact true-axis model
	cost use the same actual filling throughout.
	
	Let $Z$ have probability law $d\sigma=dA/\pi$ on $\D$, and set
	$X=2\operatorname{Re}Z$. The vertical slice at $\operatorname{Re}Z=x/2$
	has length $2\sqrt{1-x^2/4}$. Rescaling the horizontal coordinate
	gives $X$ the density $\sqrt{4-x^2}/(2\pi)$ on $[-2,2]$, zero elsewhere.
	Equivalently the angle has probability measure $dt/(2\pi)$ and
	the independent radius has measure $2r\,dr$.
	For the normalized projection $u=X/2$, put
	$g(u)=\max\{u_+^m,(1-a/m)^8u_+^{m-8}\}$.
	Then $M_m=\E[(X/2)_+^m]$, and the extra reference cost is
	\begin{equation}\label{eq:semicircle-test}
		D_f=m^{3/2}\E[g(X/2)-(X/2)_+^m]
		=\frac{m^{3/2}}{2\pi}\int_0^2
		\bigl(g(x/2)-(x/2)^m\bigr)\sqrt{4-x^2}\,dx .
	\end{equation}
	This uses the dimensionless density in \eqref{eq:signed-criterion}.
	Equivalently it is the extra unnormalized cost divided by
	$T_m=2^m m^{-3/2}$. Only $X>0$ contributes.
	With $s=m(1-x/2)$ the weighted measure $m^{3/2}\,d\mathbb P_X$
	on this half becomes $(2/\pi)\sqrt{2s-s^2/m}\,ds$ on $[0,m]$.
	The integrand difference is supported on $s\geq a$. Bounding
	$\sqrt{2s-s^2/m}$ by $\sqrt{2s}$ and integrating an exponential tail gives
	\[
	D_f\leq\frac{16\sqrt2}{\pi m}(1-a/m)^{m-1}
	\left(\frac{\sqrt a}{\sigma^2}
	+\frac1{\sqrt a\,\sigma^3}\right),
	\qquad \sigma=\frac{m-8}{m-a}.
	\]
	Indeed the two terms of $g$ cross at $s=a$; on $s\geq a$ use
	$(1-a/m)^8-(1-s/m)^8\leq8(s-a)(1-a/m)^7/m$,
	$(1-s/m)^{m-8}\leq(1-a/m)^{m-8}e^{-\sigma(s-a)}$,
	and $\sqrt s\leq\sqrt a+(s-a)/(2\sqrt a)$.
	For $m\geq1500$, $6\leq a\leq m/5$ we have
	$\sigma>0.99$, $(1-a/m)^{m-1}e^a\leq e^{a/m}<5/4$ and
	$16\sqrt2/\pi<8$.  Hence the coefficient after dividing by
	$e^{-a}\sqrt a/m$ is below
	$10((0.99)^{-2}+(6(0.99)^3)^{-1})<12$.
	Therefore
	\begin{equation}\label{eq:model-error}
		D_f\leq12e^{-a}\sqrt a/m.
	\end{equation}
	Lemma~\ref{lem:two-power-retention} gives a saving greater than
	$\Delta^3/32$ for the same test in the same direction, after
	multiplication by $m^{3/2}$.
	For the untruncated choice $a=3/y$, equation \eqref{eq:new-a} gives
	$\Delta L=9/a$ and $m=\Delta e^L/2$, with $L>9$.  Therefore
	\[
	\frac{12e^{-a}\sqrt a}{m\Delta^3}
	=\frac{24}{9^4}
	(a^4\sqrt a\,e^{-a})(L^4e^{-L}).
	\]
	For $a\geq6$, $\sqrt a\leq a/\sqrt6$ and $a^5e^{-a}$ is
	decreasing.  Also $L^4e^{-L}$ is decreasing for $L\geq9$.
	The displayed ratio is consequently at most
	$24\cdot6^4\sqrt6\,e^{-15}$.  The Taylor sum for $e^3$
	through degree eight exceeds $20$, and $\sqrt6<5/2$; hence
	this last quantity is below $243/10000<1/40$.
	Combining this with the central estimate gives
	$D_f+\text{central}<13\Delta^3/500<\Delta^3/32$.
	
	Outside the joint tiny-deficit regime, the exclusion
	\eqref{eq:tiny-cutoff} gives
	$\Delta>0.03/[m\log(100m)]$ for every remaining endpoint.
	The central and truncated costs admit a short common estimate.
	For $m\geq1500$,
	\[
	\log(100m)\leq\frac m{100}+13\leq\frac m{20},
	\qquad \Delta\geq\frac1{2m^2}.
	\]
	For the logarithm use $10000<2^{14}$, $\log2\leq1$, and
	$\log(m/100)\leq m/100-1$.
	The sequence $m^8(23/25)^m$ decreases, since
	$((m+1)/m)^8(23/25)\leq(1501/1500)^8(23/25)<1$.
	At $m=1500$, use $(23/25)^{10}<1/2$ and $1500<2^{11}$
	to obtain $m^8(23/25)^m<1/8000$.  As $\sqrt m\leq m$,
	\[
	m\sqrt m\,(23/25)^m<\frac1{8000m^6}
	\leq\frac{\Delta^3}{1000}.
	\]
	For the truncated choice $a=m/5$, the model estimate is bounded by
	this same quantity: $e^{-1/5}\leq23/25$ and
	$12\sqrt{m/5}/m\leq m\sqrt m$.
	Thus both the truncated model cost and central allowance are below
	$\Delta^3/1000$, and their sum is below $\Delta^3/64$.
	
	The central allowance applies in either phase branch.  On
	$|w|<1.84$, one has $S(w)\leq(0.92)^m$, so its contribution in
	$T_m$ units is at most $m^{3/2}(0.92)^m$.
	Outside the tiny-deficit joint regime, the lower bound for $\Delta$
	in the preceding paragraph holds for every endpoint.  The same
	rational induction therefore proves
	$m^{3/2}(0.92)^m<\Delta^3/1000$ independently of which value of
	$a$ is active.  In the true-axis branch the unscaled error
	$(0.92)^m$ is below its free model margin $M_m/m$: the largest
	binomial coefficient in \eqref{eq:reference-moment} is at least
	$2^m/(m+1)$, so $M_m/m\geq1/[m(m+1)(m+2)]$, whereas
	$m^8(23/25)^m<1/8000$.
	
	\begin{lemma}[Reduction to a tiny first endpoint]\label{lem:tiny-first}
		For an actual nonslit weak Catalan candidate, the same selected first
		endpoint satisfies
		\[
		0< m d_1 < \frac{3}{100m\log(100m)}.
		\]
		Every unit direction has the global cap $2-d_1$ and the actual
		disk saving $>(md_1)^3/30$ in the normalization used above.
	\end{lemma}
	\begin{proof}
		If both scaled deficits were at least the displayed threshold,
		the aligned model estimate and the central estimate would make
		each aligned endpoint budget strictly negative.  The true-axis
		reference margin also pays this allowance, independently of the
		deficit.  The full-disk comparison counts the center only once,
		so the candidate lower bound is impossible.  At least one
		endpoint is therefore tiny.  Global radius maximality gives
		$d_1\leq d_2$, so it is the first endpoint that is tiny.  Its
		global cap holds in every direction, and the actual capped-power
		saving uses only that cap and the same radius.  This reduction
		and the reference-tail estimate use the same endpoint. The preceding
		antipodal comparison treats the remaining tiny range.
	\end{proof}
	
	\begin{proof}[Proof of Theorem~\ref{thm:large}]
		Suppose first that an actual nonslit map is a weak contender:
		$|B_n(F)|\geq\Cat_{k-1}$.  The boundary localization gives two
		actual assigned ends with $0<nd_i<1$, their whole-hull strips and
		whole-boundary projection caps, and the joint strict scalar cutoff.
		Lemma~\ref{lem:cutoff} then gives $d_i<1/10000$ and
		$\Delta_iL_i<3/2$ on both sides.  The refined root estimate gives
		$mW_i\leq y_i<1/2$, and the preceding angle argument gives the
		cap for each chosen true axis or nearest phase ray.  The energy
		multiplier in Section~\ref{sec:energy-cert} applies to this actual
		map, so the signed comparison \eqref{eq:signed-criterion} holds.
		
		If the minimum deficit satisfies \eqref{eq:tiny-cutoff}, the two
		opposite capped rays and the central estimate above give
		$\int_{\D}S(H)\,d\sigma<2M_m$, contradicting the candidate
		inequality.  Otherwise apply the outer-region comparisons to each
		assigned side.  In the aligned branch choose the nearest phase ray
		and the convex test $g(u)=\max\{u_+^m,(1-a/m)^8u_+^{m-8}\}$.
		The root strip gives the displayed aligned tilt bound, so the
		pointwise signed integrand is below $g$ on $|H|\geq1.84$.
		The central remainder is at most $(0.92)^m$.
		For this direction, Lemma~\ref{lem:two-power-retention} saves more
		than $\Delta_i^3/(32m^{3/2})$ from the reference two-power integral.
		That reference integral is $M_m+D_f/m^{3/2}$. Even charging the
		whole central remainder to this side yields an upper bound strictly below
		\[
		M_m+\frac{D_f-\Delta_i^3/32}{m^{3/2}}+(0.92)^m .
		\]
		By the preceding estimates $D_f+m^{3/2}(0.92)^m
		<13\Delta_i^3/500<\Delta_i^3/32$. This leaves the cost strictly
		below $M_m$.
		
		In the other branch use the true endpoint axis and the convex test
		$g_0(u)=\lambda u_+^{m-8}$.  The true-axis tilt bound and
		\eqref{eq:true-axis-corrected} give strict pointwise domination
		on the outer region.  The uncut model integral is
		$(1-1/m)M_m$; the cap can only decrease it, and the central
		remainder is less than $M_m/m$.  This side also costs strictly
		below $M_m$.  The direction is fixed on every source circle in each
		convex-order application.  Summing both assigned sides contradicts
		$\int S(H)\geq2M_m$.  Thus every nonslit actual map has strict
		inequality, while a zero endpoint deficit is already a slit by the
		area rigidity following \eqref{eq:two-point}.
		
		Finally $F_a(z)=z+a/z$, $|a|=1$, has inverse expansion
		$G_a(w)=w-\sum_{k\geq1}\Cat_{k-1}a^k w^{-(2k-1)}$.
		It attains the bound in every listed order.  This proves both the
		upper bound and the full equality classification.
	\end{proof}
	\subsection{The phase-energy inequality}\label{sec:energy-cert}
	The energy estimate controls the phase selected by the inverse
	coefficient.  Its proof uses a drift identity, one signed flux, and
	integration with a radial cutoff. We choose a single capped linear
	function to obtain a constant uniform in $k$.
	
	\begin{lemma}\label{lem:actual-phase-energy}
		For every centered exterior univalent $F$, integer $k\geq751$, and
		$|\beta|=1$, put $H(z)=z+g(\bar z)$ and
		$U(z)=\operatorname{Im}(\beta H(z)^k)$.  Then
		\[
		k^2\int_\D \frac{U^2}{|H|^2}\,d\sigma
		\leq\frac73\int_\D |\nabla U|^2\,d\sigma.
		\]
		The quotient is extended by zero at $H=0$.
	\end{lemma}
	\begin{proof}
		Write $y=|H|$ and $c=g'(\bar z)$.  Real differentiation gives
		\begin{equation}\label{eq:energy-drift}
			y^2\Delta U+4k(k-1)U=4(k-1)H\cdot\nabla U.
		\end{equation}
		This identity extends across $H=0$. In fact,
		$\Delta U=4k(k-1)\operatorname{Im}(\beta H^{k-2}c)$;
		the terms containing $g''$ cancel.  
		
		For $y>0$ put $u=U/y$, $p=|\nabla U|$, $r=|z|$, $t=k(1-r)/y$
		and $b=k/[4(k-1)]$. Choose
		\[
		\chi(t)=\min\{12t,9/5\},\qquad f(t)=\chi(t)/6,\qquad
		J=\chi(t)U\nabla U+4(k-1)e_r f(t)\frac{U^2}{y}.
		\]
		Combine the two divergences first. The coefficient of $uU_r$
		from differentiating $\chi$ and the radial flux is
		$8(k-1)f-k\chi'$. For the radial term,
		with $h=f+tf'$, the product rule reads
		\begin{equation}\label{eq:energy-radial-identity}
			\operatorname{div}\left(e_r f(t)U^2/y\right)
			=\{(y/r)f-hy_r-kf'\}u^2+2fuU_r.
		\end{equation}
		Use \eqref{eq:energy-drift} together with
		\[
		\nabla t=-\frac{ke_r+t\nabla y}{y},\quad
		|u||\nabla y|\leq2|u|+p/k,\quad
		(y/r)f\leq2f.
		\]
		For the gradient bound, $|\nabla y|\leq1+|c|$ and $|u|\leq y^{k-1}$.
		Differentiating the imaginary part gives
		$p=k|\beta H^{k-1}-\overline{\beta H^{k-1}c}|
		\geq k y^{k-1}\bigl||c|-1\bigr|$.
		Define
		\begin{align*}
			A_k&=\chi+f'-(2f+2|h|)/k,\\
			D_k&=(k\chi+t|\chi'|)/(4(k-1)),\\
			B_k&=\chi+|2f-b\chi'|+t|\chi'|/(2(k-1))+|h|/k .
		\end{align*}
		Indeed, the drift and radial identities give
		\begin{align*}
			b\operatorname{div}J+k^2(\chi+f')u^2
			&=b\chi p^2+k\chi u(H/y)\cdot\nabla U
			+k(2f-b\chi')uU_r\\
			&\quad-bt\chi' u\,\nabla y\cdot\nabla U
			+k\bigl((y/r)f-hy_r\bigr)u^2 .
		\end{align*}
		Using the preceding bounds in this equality yields
		\[
		k^2A_k(t)u^2+b\operatorname{div}J
		\leq D_k(t)p^2+kB_k(t)|u|p.
		\]
		The scalar multiplier bound below implies its homogeneous form
		also when $p=0$, hence
		\begin{equation}\label{eq:energy-pointwise-flux}
			k^2u^2+b\operatorname{div}J\leq\frac73 p^2
			\quad\hbox{almost everywhere on }\D.
		\end{equation}
		
		At the single multiplier corner the differential of $t$ is zero almost
		everywhere on that level set.  The Lipschitz multiplier compositions
		therefore have zero differential there; either adjacent slope gives
		the same product rule.  No nullity assumption on the level is needed.
		At a filling zero, $|J|=O(y^{2k-1})$ and $y=O(|z-z_0|)$,
		so its zero extension has derivative zero.
		
		For local Lipschitz continuity through the zero set, use the following
		elementary estimate.  If $u,v$ are Lipschitz, with constants
		$K_u,K_v$, and $|u|\leq C|v|$, the zero-extended quotient $u^2/v$
		is Lipschitz with constant $2CK_u+C^2K_v$.  When $|b|\leq|a|$
		and $ab\ne0$, use
		\[
		\frac{x^2}{a}-\frac{y^2}{b}
		=\frac{(x-y)(x+y)}a+\frac{y^2(b-a)}{ab};
		\]
		if $b=0$, then $y=0$.  Apply this to the actual phase field and
		$v=|H|$, using $|U|\leq2^{k-1}|H|$.
		Near a noncentral filling zero the multipliers have constant tail
		values wherever $H\ne0$, and the amplitudes vanish at the other
		zeros.  Thus $J$ is locally Lipschitz on the punctured source disk.
		
		The area theorem gives $g'\in L^2(\D)$, hence $p^2,p\in L^1(\D)$;
		also $u^2\leq2^{2k-2}$.  Bounded multipliers give
		$|J|\leq C_k y^{2k-1}(1+|g'|)$, while their linear vanishing
		at $t=0$ gives $|J|\leq C_k(1-r)y^{2k-2}(1+|g'|)$.
		Using $y\leq2r$ and $y\leq2$, respectively, obtain
		\begin{equation}\label{eq:energy-boundary-majorant}
			|J(z)|\leq C_k\min\{r,1-r\}\bigl(1+|g'(\bar z)|\bigr).
		\end{equation}
		
		Let $\eta_\varepsilon$ be the product of two unit linear ramps,
		zero for $r\leq\varepsilon$ or $r\geq1-\varepsilon$, one for
		$2\varepsilon\leq r\leq1-2\varepsilon$, with Lipschitz constant
		at most $2/\varepsilon$.  The field $\eta_\varepsilon J$ is
		globally Lipschitz with compact support.  Absolute continuity along
		horizontal and vertical lines, the one-dimensional fundamental
		theorem of calculus and Fubini give
		\[
		\int_\D\eta_\varepsilon\operatorname{div}J\,d\sigma
		=-\int_\D\nabla\eta_\varepsilon\cdot J\,d\sigma.
		\]
		Both integrands are integrable.  By
		\eqref{eq:energy-boundary-majorant},
		$|\nabla\eta_\varepsilon\cdot J|
		\leq4C_k(1+|g'(\bar z)|)$.
		The derivative of the cutoff is eventually zero at each noncentral
		interior point, so dominated convergence removes the boundary term.
		Multiply \eqref{eq:energy-pointwise-flux} by the cutoff and pass to
		the limit in the $L^1$ densities $u^2,p^2$.  No prior integrability
		assumption on the unweighted divergence is required.
	\end{proof}
	
	\paragraph{The scalar multiplier bound.}
	On the increasing part write $x=12t\in[0,9/5]$. Use
	$a=A_k-1$, $\mathcal B=B_k$, $\mathcal D=D_k$ within this calculation only.
	Here $2f\leq3/5<3k/(k-1)$, so the absolute value in $B_k$ has a fixed
	sign. The exact coefficients are
	\[
	\begin{aligned}
		a&=1+(1-1/k)x,\\
		\mathcal B&=3k/(k-1)+\bigl(2/3+1/(2(k-1))+1/(3k)\bigr)x,\\
		\mathcal D&=\bigl(1+2/(k-1)\bigr)x/4 .
	\end{aligned}
	\]
	On the constant part $\chi=9/5$, $f=3/10$, and these become
	\[
	a=4/5-6/(5k),\qquad \mathcal B=12/5+3/(10k),\qquad
	\mathcal D=9k/(20(k-1)).
	\]
	Let $(a_0,b_0,d_0)$ be the limits of $(a,\mathcal B,\mathcal D)$
	as $k\to\infty$ at fixed $t$.
	To see the reason for the cap $9/5$, allow an arbitrary plateau
	$\chi=c>1$ while retaining $f=\chi/6$. On the constant part the
	limiting cost satisfies
	\[
	\frac c4+\frac{4c^2}{9(c-1)}-\frac94
	=\frac{(5c-9)^2}{36(c-1)} .
	\]
	The cap $c=9/5$ therefore gives precisely the cost $9/4$ at the
	origin of the increasing part as well as on the plateau.
	Inside the increasing part one has
	\[
	4a_0(9/4-d_0)-b_0^2
	=(1+x)(9-x)-(3+2x/3)^2
	=x(4-13x/9)\geq0 .
	\]
	On the plateau $(a_0,b_0,d_0)=(4/5,12/5,9/20)$. Thus $d_0+b_0^2/(4a_0)\leq9/4$ throughout.
	
	For $k\geq751$ the two explicit formulas above give
	\[
	a\geq(99/100)a_0,\qquad
	0\leq \mathcal B\leq(101/100)b_0,\qquad
	0\leq \mathcal D\leq(101/100)d_0 .
	\]
	For example, on the increasing part the excess in $\mathcal B$ is at most
	$(3+5x/6)/(k-1)\leq b_0/100$; on the plateau the relative loss
	in $a$ is $3/(2k)$. The remaining bounds follow directly from
	$2/(k-1)<1/100$. Completing the square now gives, for all $p$ and $u$
	with $p\geq0$,
	\[
	kp \mathcal B |u|-k^2a u^2+\mathcal D p^2
	\leq \left(\mathcal D+\frac{\mathcal B^2}{4a}\right)p^2
	\leq \frac{101^2}{100\cdot99}\frac94\,p^2
	< \frac73\,p^2\quad (p>0).
	\]
	For $p=0$ the left side is nonpositive. Either one-sided expression
	at the corner has this bound, as required in the product rule.
	
	\section{Statement and boundary trace}
	Let $\Sigma_0$ be the class of injective holomorphic functions on $|z|>1$ with expansion
	\begin{equation}\label{eq:F}
		F(z)=z+\sum_{j\ge1}b_jz^{-j}.
	\end{equation}
	The absence of a constant term fixes the center. For $p>0$, let
	\[
	R_p(F)=\frac1{2\pi}\int_{-\pi}^{\pi}|F^*(e^{it})|^p\,dt,
	\qquad S(p)=\frac{\Gamma(p+1)}{\Gamma(1+p/2)^2},
	\]
	where $F^*$ is the almost-everywhere radial limit, constructed below. Write
	$D=1-|b_1|^2$. The area theorem gives $D\ge0$.
	
	Baernstein, Laugesen, and Pritsker \cite[Section 6, pp.~78--79]{BLP}
	already proved $R_2(F)\le2$ in this exterior normalization and showed
	that equality is possible exactly for a rotation of $[-2,2]$. Their
	argument uses Parseval and the area theorem. We retain the unused
	area budget to obtain a quantitative second-moment deficit. For the
	next three even exponents, the same area budget controls a finite
	weighted sum after the signed-area identity below; low-order Grunsky
	inequalities handle the remaining coefficient combinations.
	
	\begin{theorem}\label{thm:moments}
		For every $F\in\Sigma_0$,
		\begin{align}
			R_2(F)&\le 2-\tfrac12D=S(2)(1-D/4),\label{eq:m2}\\
			R_4(F)&\le 6-3D=S(4)(1-D/2),\label{eq:m4}\\
			R_6(F)&\le 20-7D-\tfrac92D^2,\label{eq:m6}\\
			R_8(F)&\le 70-\tfrac{51}{5}D.\label{eq:m8}
		\end{align}
		For each $p\in\{2,4,6,8\}$, $R_p(F)=S(p)$ if and only if
		$F(z)=z+a/z$ for some $|a|=1$. Thus the four maximal values
		$S(2)=2$, $S(4)=6$, $S(6)=20$, and $S(8)=70$ are sharp.
	\end{theorem}
	
	The deficit coefficients in the four inequalities are explicit constants
	from the proof. The sharpness statement concerns the four maximal values.
	
	\begin{lemma}\label{lem:trace}
		The radial limit $F^*(e^{it})=\lim_{r\downarrow1}F(re^{it})$ exists for almost
		every $t$, and for every fixed $p>0$ its $p$th power is integrable. Moreover,
		\[
		R_p(F)=\lim_{r\downarrow1}\frac1{2\pi}\int_{-\pi}^{\pi}
		|F(re^{it})|^p\,dt.
		\]
	\end{lemma}
	\begin{proof}
		Set $g(u)=F(1/u)-1/u=\sum_{j\ge1}b_ju^j$. The area theorem gives
		$\sum_{j\ge1}j|b_j|^2\le1$. Parseval on concentric disks and monotone
		convergence therefore give
		\[
		\int_{\D}|g'(u)|^2\,dA(u)=\pi\sum_{j\ge1}j|b_j|^2\le\pi.
		\]
		Fubini implies that for almost every $t$ the radial integral of
		$|g'(\rho e^{it})|^2$ on $1/2<\rho<1$ is finite. Since $\rho\ge1/2$,
		Cauchy--Schwarz shows that $g(\rho e^{it})$ is Cauchy as $\rho\uparrow1$.
		This constructs the asserted trace.
		
		For completeness, the omitted-point bound makes the boundary passage
		dominated. If $w$ is omitted by a centered exterior map $H$, the function
		$u\mapsto 1/(H(1/u)-w)=u+wu^2+\cdots$ is schlicht on $\D$.
		The second-coefficient theorem gives $|w|\le2$. For each $r>1$, set
		$F_r(z)=F(rz)/r$. This is a centered normalized univalent map on
		$|z|>1/r$. In particular, for $|z|=1$ the point $F_r(z)$ is omitted
		by the restriction of $F_r$ to $|z|>1$, by injectivity on the larger
		domain. Hence $|F_r(e^{it})|\le2$, or
		$|F(re^{it})|\le2r$. This uniform bound for $r$ in a fixed interval
		$(1,r_0)$ gives dominated convergence on the circles as $r\downarrow1$.
	\end{proof}
	
	\section{Weighted area and the first four even moments}
	We use the classical area theorem and strong Grunsky inequality in their
	usual exterior normalization \cite[Chapters 1 and 3]{Pom}. The complex
	Grunsky coefficients used below are $c_{22}$, $c_{23}$, and $c_{33}$;
	their Laurent expressions are displayed in the proof.
	
	\begin{lemma}[Finite radial area inequality]\label{lem:area}
		Put $\phi(u)=1+\sum_{j\ge1}b_ju^{j+1}$ and
		$c_d=[u^d]\phi(u)^m$. Then, for each positive integer $m$,
		\begin{equation}\label{eq:rad-area}
			R_{2m}(F)\le \sum_{d=0}^{2m-1}\left(2-\frac dm\right)|c_d|^2.
		\end{equation}
	\end{lemma}
	\begin{proof}
		First replace $F$ by $F_r(z)=F(rz)/r$, $r>1$. This map is univalent on
		a neighborhood of the closed exterior disk. On the unit circle,
		$F_r(z)^m=\sum_{d\ge0}c_d^{(r)}z^{m-d}$. Parseval gives its mean squared
		modulus as $\sum_d|c_d^{(r)}|^2$. Let $K_r$ be the bounded region
		enclosed by the smooth Jordan curve $F_r(\partial\D)$. Apply Green's
		formula on $K_r$ to the analytic map $w\mapsto w^m$. The oriented image
		contour, counted with its multiplicity, has signed area
		$m^2\int_{K_r}|w|^{2m-2}\,dA(w)\ge0$. Fourier orthogonality writes the
		same signed area as
		$\pi\sum_d(m-d)|c_d^{(r)}|^2$. Adding this nonnegative sum to $m$ times
		the Parseval sum gives
		\[
		m\sum_{d\ge0}|c_d^{(r)}|^2
		\le\sum_{d\ge0}(2m-d)|c_d^{(r)}|^2
		\le\sum_{d<2m}(2m-d)|c_d^{(r)}|^2.
		\]
		The last inequality discards only nonpositive terms. Divide by $m$,
		then let $r\downarrow1$ using Lemma~\ref{lem:trace}; the sum on the
		right is finite and each coefficient converges.
	\end{proof}
	
	\begin{proof}[Proof of Theorem~\ref{thm:moments}]
		Write $s=|b_1|^2$, so $0\le s\le1$ and $D=1-s$. Parseval and the area
		theorem yield
		\[
		R_2(F)=1+\sum_{j\ge1}|b_j|^2
		\le1+s+\frac12\sum_{j\ge2}j|b_j|^2
		\le1+s+D/2=2-D/2.
		\]
		At $m=2$, equation~\eqref{eq:rad-area} expands to
		\[
		R_4(F)\le2+4s+2|b_2|^2\le2+4s+D=6-3D.
		\]
		
		For $m=3$, the coefficients $c_0,\ldots,c_5$ of $\phi^3$ are
		$1,0,3b_1,3b_2,3(b_3+b_1^2),3(b_4+2b_1b_2)$. Thus
		\begin{equation}\label{eq:m6finite}
			R_6(F)\le2+12s+9|b_2|^2+6|b_3+b_1^2|^2
			+3|b_4+2b_1b_2|^2.
		\end{equation}
		Put $x=|b_1|$, $a=|b_2|$, and $b=|b_4|$. Then $0\le x\le1$,
		$|b_4+2b_1b_2|\le b+2xa$, and the square completion is
		\[
		(5+6x^2)(2a^2+4b^2)-9a^2-3(2xa+b)^2
		=(a-6xb)^2+(17-12x^2)b^2\ge0.
		\]
		The area theorem therefore bounds the two terms containing $b_2,b_4$
		in \eqref{eq:m6finite} by $(5+6x^2)D$.
		
		For the remaining term, define the exterior Grunsky coefficients by
		\[
		-\log\frac{F(z)-F(w)}{z-w}
		=\sum_{j,k\ge1}c_{jk}z^{-j}w^{-k}.
		\]
		Expansion gives $c_{22}=b_3+b_1^2/2$. The strong Grunsky inequality
		applied to the vector supported at index two says
		$\sum_{j\ge1}j|c_{2j}|^2\le1/2$; hence $|c_{22}|\le1/2$.
		Consequently $|b_3+b_1^2|\le(1+s)/2$. Combining the estimates gives
		\[
		R_6(F)\le2+12s+(5+6s)D+\frac32(1+s)^2
		=20-7D-\frac92D^2.
		\]
		For $m=4$ put
		\[
		u=b_4+b_1b_2,\quad v=b_1b_2,\quad g=b_3+b_1^2,
		\quad A=4b_3+6b_1^2,
		\]
		\[
		w=b_5+3b_1b_3+\tfrac32b_2^2+b_1^3,
		\qquad y=b_6+3b_1b_4+3gb_2.
		\]
		Expanding the first eight coefficients in \eqref{eq:rad-area} gives
		\begin{equation}\label{eq:m8finite}
			R_8(F)\le2+24s+20|b_2|^2+|A|^2
			+12|u+2v|^2+8|w|^2+4|y|^2.
		\end{equation}
		The Grunsky coefficients are $c_{23}=u$ and
		$c_{33}=b_5+b_1b_3+b_2^2+b_1^3/3$.
		Applying the strong inequality to vectors supported at indices two
		and three gives
		\begin{equation}\label{eq:grunsky8}
			|b_2|^2+2|c_{22}|^2+3|u|^2\le\tfrac12,
			\qquad |c_{33}|\le\tfrac13.
		\end{equation}
		
		First,
		\[
		|A|^2=48|c_{22}|^2+24s^2-32|b_3|^2,
		\qquad
		12|u+2v|^2-72|u|^2\le\tfrac{288}{5}|v|^2.
		\]
		The second inequality follows on replacing $|u+2v|$ by
		$|u|+2|v|$ and completing the square:
		$\frac{288}{5}|v|^2-12(|u|+2|v|)^2+72|u|^2
		=\frac{12}{5}(5|u|-2|v|)^2$.
		Since $|v|^2=s|b_2|^2\le sD/2$, \eqref{eq:grunsky8} yields
		\begin{equation}\label{eq:firstblock}
			20|b_2|^2+|A|^2+12|u+2v|^2
			\le12+24s^2+\tfrac{144}{5}sD.
		\end{equation}
		
		Next, $w=c_{33}+(2b_1b_3+\frac23b_1^3)+\frac12b_2^2$.
		Expansion of $|b_3+b_1^2/2|^2\le1/4$ gives
		\[
		|2b_1b_3+\tfrac23b_1^3|^2
		\le\tfrac43s|b_3|^2+\tfrac23s-\tfrac29s^3
		\le(\tfrac23+s|b_3|^2)^2.
		\]
		For the last inequality the difference of the right and left sides
		is $\frac29(1-s)^2(s+2)+(s|b_3|^2)^2$.
		Using \eqref{eq:grunsky8} and
		$3|b_3|^2+2|b_2|^2\le D$ gives
		$|w|\le1+s|b_3|^2+|b_2|^2/2\le1+D/3$.
		As $0\le D\le1$,
		\begin{equation}\label{eq:wblock}
			8|w|^2\le8+\tfrac{16}{3}D+\tfrac89D^2
			=8+\bigl(\tfrac{56}{9}-\tfrac89s\bigr)D.
		\end{equation}
		
		Finally, weighted Cauchy--Schwarz and the area theorem give
		\[
		|y|^2\le
		\left(\tfrac16+\tfrac94s+\tfrac92|g|^2\right)
		(6|b_6|^2+4|b_4|^2+2|b_2|^2).
		\]
		The last factor is at most $D$, and the diagonal Grunsky estimate
		above gives $|g|\le(1+s)/2$. Hence
		\begin{equation}\label{eq:yblock}
			4|y|^2\le
			\left(\tfrac23+9s+\tfrac92(1+s)^2\right)D.
		\end{equation}
		Combining \eqref{eq:firstblock}--\eqref{eq:yblock} with
		\eqref{eq:m8finite} gives $70-R_8(F)\ge D\widetilde P(s)$, where
		\[
		\widetilde P(s)=48+24s-\tfrac{144}{5}s-\tfrac{56}{9}
		-\tfrac23-9s-\tfrac92(1+s)^2+\tfrac89s.
		\]
		Direct expansion gives
		\[
		\widetilde P(s)-\frac{51}{5}
		=\frac{(1-s)(405s+2377)}{90}\ge0\qquad(0\le s\le1),
		\]
		which proves \eqref{eq:m8}.
		
		The values $S(2)=2$, $S(4)=6$, $S(6)=20$, and $S(8)=70$ follow from the gamma
		recurrence. Equality with $S(p)$ in any displayed bound forces $D=0$.
		Then the area theorem forces $b_j=0$ for every $j\ge2$, so
		$F(z)=z+b_1/z$ with $|b_1|=1$. Conversely, that map has the rotated
		segment $[-2,2]$ as its boundary image, and its even moments are
		$\binom{2m}{m}$ for $m=1,2,3,4$ by the elementary cosine integral.
	\end{proof}
	
	\section{Every real exponent from four to eight}\label{sec:continuous-moments}
	The four endpoint estimates also control every exponent between four and
	eight.  We use one published comparison in addition to the ingredients of
	Theorem~\ref{thm:moments}.  For a centered capacity-one continuum, the
	Baernstein--Laugesen--Pritsker projection theorem
	\cite[Theorem~2]{BLP} says that every real projection of its equilibrium
	measure is dominated, for convex tests, by the arcsine distribution on
	$[-2,2]$.  Applied to the omitted continuum of $F$, rotated by a unit
	complex number $c$, and using the harmonic-measure boundary
	correspondence, this gives
	\begin{equation}\label{eq:projection-order}
		\E\,\psi\!\left(\frac{\operatorname{Re}(cF^*)}{2}\right)
		\le \E\,\psi(\cos t)
	\end{equation}
	for every convex $\psi:\mathbb R\to\mathbb R$. Filling bounded complementary
	components changes neither the exterior map nor its harmonic measure;
	its capacity and conformal centroid are respectively one and zero by
	the normalization in \eqref{eq:F}.  Thus the cited theorem has precisely
	the hypotheses required here.  The boundary passage in
	Lemma~\ref{lem:trace} gives the same comparison for the actual radial
	trace.
	
	\begin{theorem}\label{thm:radial-continuous}
		For every $F\in\Sigma_0$, every real $4\le p\le8$, and
		$D=1-|b_1|^2$,
		\begin{equation}\label{eq:continuous-deficit}
			R_p(F)\le S(p)(1-D/80).
		\end{equation}
		For each fixed $p$ in this interval,
		$R_p(F)=S(p)$ if and only if $F(z)=z+a/z$ on $|z|>1$ for some $|a|=1$.
	\end{theorem}
	
	\begin{proof}
		Normalize $T=|F^*|/2$, and write
		$C_q=\E|\cos t|^q=2^{-q}S(q)$.  Theorem~\ref{thm:moments} gives
		\begin{equation}\label{eq:three-normalized-deficits}
			C_j-\E T^j\ge d_jD,\qquad
			(d_4,d_6,d_8)=\left(\frac3{16},\frac7{64},\frac{51}{1280}\right).
		\end{equation}
		For $j=6$, division of $20-R_6\ge7D+\frac92D^2$ by $2^6$
		gives a further nonnegative term; the other two values follow by
		dividing the fourth- and eighth-moment deficits by $2^4$ and $2^8$.
		
		Choose the phase $c$ with $|c|=1$ so that the first Fourier mode of
		$X=\operatorname{Re}(cF^*)/2$ has amplitude $(1+|b_1|)/2$.
		Parseval, first on circles and then for the actual trace, yields
		\[
		\E X^2\ge\frac{(1+|b_1|)^2}{8},\qquad
		\frac12-\E X^2\le\frac{3D}{8}.
		\]
		The last inequality is the elementary identity
		$3(1-x^2)-(4-(1+x)^2)=2x(1-x)\ge0$ for $x=|b_1|\in[0,1]$.
		
		For $a=4$ or $6$, put $p=a+2\theta$ with $0\le\theta\le1$, and let
		\[
		h(u)=u^p-\frac{p(1-\theta)}{a}u^a
		-\frac{p\theta}{a+2}u^{a+2},\qquad
		L_a=\frac1{6(a-2)}+\frac4{7(a-2)^2}.
		\]
		Weighted arithmetic--geometric mean gives $h'(u)\le0$ on $[0,1]$.
		The second derivative satisfies
		\begin{equation}\label{eq:h-curvature}
			h''(u)\ge-2p\theta(1-\theta)L_a\qquad(0\le u\le1).
		\end{equation}
		Here is the short calculus check of the constant.  If $u=e^{-v}$,
		the secant error in $h''/p$ is that of
		$f(k)=(k-1)e^{-(k-2)v}$ between $a$ and $a+2$.
		With $w=(k-2)v$,
		\[
		f''(k)=e^{-w}\left(\frac{w^2-2w}{k-2}
		+\frac{w^2}{(k-2)^2}\right)
		\le L_a,
		\]
		because the two maxima for $w\ge0$ are respectively
		$(2+2\sqrt2)e^{-2-\sqrt2}<1/6$ and $4e^{-2}<4/7$.
		For the first comparison, \(7/5<\sqrt2<10/7\) and
		\(\sum_{n=0}^7(17/5)^n/n!>204/7\) give
		\(e^{2+\sqrt2}>12(1+\sqrt2)\).
		For the second, the Taylor sum through degree four at \(2\) equals
		\(7\), and every remaining term is positive.
		The secant error over a length-two interval is at most
		$2\theta(1-\theta)L_a$, proving \eqref{eq:h-curvature}.
		As $h'(0)=0$, the even function
		$u\mapsto h(|u|)+p\theta(1-\theta)L_a u^2$ is convex on $[-1,1]$.
		Continue it by its endpoint tangent lines to obtain a convex test
		on $\mathbb R$.
		
		Since $|X|\le T\le1$ and $h$ is decreasing, $h(T)\le h(|X|)$.
		Apply \eqref{eq:projection-order} to the preceding convex test and
		expand both averages.  Inserting \eqref{eq:three-normalized-deficits}
		and the projection second-moment bound gives
		\begin{equation}\label{eq:interpolation-bracket}
			C_p-\E T^p\ge pD B_a(\theta),\qquad
			B_a(\theta)=(1-\theta)A+\theta B-L\theta(1-\theta),
		\end{equation}
		where $A=d_a/a$, $B=d_{a+2}/(a+2)$ and $L=3L_a/8$.
		The numerical data are
		\[
		\begin{array}{c|ccc}
			a&A&B&L\\ \hline
			4&3/64&7/384&19/224\\
			6&7/384&51/10240&13/448
		\end{array}
		\]
		and need only four fraction comparisons.  Indeed
		\[
		4-3\theta-9\theta(1-\theta)=(3\theta-2)^2\ge0
		\]
		implies
		\[
		B_a(\theta)\ge
		(1-\theta)(A-4L/9)+\theta(B-L/9).
		\]
		For $a=4$ the two endpoint values are $37/4032$ and $71/8064$;
		for $a=6$ they are $43/8064$ and $1133/645120$.
		Each is strictly greater than $1/640$.  Hence
		$C_p-\E T^p\ge pD/640\ge D/160$.
		Finally $C_p\le C_2=1/2$ for $p\ge2$, so
		$D/160\ge C_pD/80$.  Multiplication by $2^p$ proves
		\eqref{eq:continuous-deficit}.
		
		If $R_p(F)=S(p)$, then \eqref{eq:continuous-deficit} and $S(p)>0$
		force $D=0$.  The area theorem gives $b_j=0$ for $j\ge2$ and
		$|b_1|=1$, so $F$ is a unit rotated slit on the entire open exterior
		disk.  Conversely the slit boundary is a rotated $2\cos t$, and its
		moment equals $S(p)$.
	\end{proof}
	\section{Zeros of polynomials with connected lemniscates}
	\label{sec:polynomial-zero-comparison}
	
	\begin{proposition}\label{prop:zero-moments}
		Let $P$ be a complex polynomial of degree $n\ge2$, with leading
		coefficient $a$. Let $c$ be the mean of its zeros, counted with
		multiplicity. Suppose that $t>0$ and that $\{|P|\le t\}$ is connected,
		and put $R=(t/|a|)^{1/n}$. For every real $4\le p\le8$ and every
		integer $0\le j<n$,
		\[
		\frac1{n-j}\sum_{P^{(j)}(z)=0}
		\left|\frac{z-c}{R}\right|^p < S(p).
		\]
		The sum counts multiplicities. For each fixed $p$ and $j$, the
		supremum of these means over all such polynomials and levels,
		allowing $n$ to vary, is $S(p)$ and is never attained.
	\end{proposition}
	
	\subsection{The normalized exterior map}
	Put $Q(z)=P(c+Rz)/(aR^n)$. This polynomial is monic, its roots sum to
	zero, and $K=\{|Q|\le1\}$ is connected. The chain rule gives
	$Q^{(j)}(z)=R^j P^{(j)}(c+Rz)/(aR^n)$, so the affine change of
	variables preserves derivative-root multiplicities and reduces the
	mean to $(n-j)^{-1}\sum_{Q^{(j)}(z)=0}|z|^p$.
	
	We construct the normalized exterior map directly. First,
	$\Omega=\{|Q|>1\}$ is connected. The polynomial is open and proper,
	and its restriction to $\Omega$ maps into the exterior unit disk.
	Each component of $\Omega$ is open and closed relative to $\Omega$.
	Its image under the restriction is nonempty, open and closed
	relative to the connected exterior unit disk. Consequently each
	component maps onto that disk and is unbounded. Since $K$ is compact,
	all these components meet the same connected region outside a
	sufficiently large circle.
	
	Choose $b\in K$ and put $K_b=K-b$. In the coordinate $u=(z-b)^{-1}$,
	the reciprocal exterior domain $U$ consists of zero and the inverses
	of points of $\C\setminus K_b$. It is open and connected, since the
	complement of $K_b$ is connected, omits zero, and contains all points
	of sufficiently large modulus. For $\epsilon>0$, use the open region
	\[
	U_\epsilon=\{u: \epsilon |u|<|1-vu|\text{ for all }v\in K_b\}.
	\]
	It is open by compactness of $K_b$. Every point of $U$ is in the
	component containing zero of some $U_\epsilon$: join it to zero
	by a path in $U$ and take the minimum of $|1-vu|$ and a bound for
	$|u|$ on the compact product of the two parameter sets.
	
	For $v\in K_b$, join $v$ to zero by a finite chain in $K_b$ of steps
	shorter than $\epsilon$. Such chains exist because the set reachable
	by short steps is both relatively open and relatively closed in
	$K_b$. For consecutive points $v_i,v_{i+1}$ of the chain, the ratio
	$(1-v_i u)/(1-v_{i+1}u)$ lies at distance less than one from $1$ on
	$U_\epsilon$. Sum its principal logarithms along the chain to get
	an analytic logarithm of $1-vu$ with value zero at $u=0$.
	On the components containing zero the resulting logarithms agree
	where both are defined. Indeed the smaller component lies in the
	larger one, and the branches agree near zero by local injectivity of
	the exponential. The identity theorem gives agreement on that
	component. Thus there is a normalized analytic logarithm on $U$ for
	every $v\in K_b$.
	
	Let $q_1,\ldots,q_n$ be the roots of $Q$, with multiplicities.
	The exponential of the mean of the preceding logarithms for
	$v=q_i-b$ is an analytic function $H$ on $U$ with
	\[
	H(0)=1,\qquad H(u)^n=\prod_i(1-(q_i-b)u),\qquad H'(0)=b.
	\]
	The last identity follows by differentiation and $\sum_i q_i=0$.
	Hence
	$B(z)=(z-b)H((z-b)^{-1})$ is analytic on $\Omega$ and $B(z)^n=Q(z)$.
	It maps $\Omega$ into the exterior unit disk and is proper there.
	For instance, a compact set of target points has compact image
	under the $n$th power map, and its inverse image under $B$ is
	closed in the compact inverse image under $Q$ of that image.
	Since $B$ is nonconstant and analytic, its image is open. Properness
	makes the image closed relative to the exterior unit disk, so it is
	onto. To show injectivity, take any exterior value $w$ and pick a
	preimage under $B$ of each of the $n$ distinct $n$th roots of $w^n$.
	These yield $n$ distinct roots of $Q(z)-w^n$, which exhaust its
	degree. Each chosen value has exactly one preimage.
	
	The inverse $F$ of $B$ therefore satisfies $Q(F(w))=w^n$ for $|w|>1$.
	As $w\to\infty$, also $F(w)\to\infty$ by this identity. With
	$u=(F(w)-b)^{-1}$ we have
	$F(w)-w=b-(H(u)-1)/u\to b-H'(0)=0$.
	The reciprocal correction thus extends holomorphically to zero
	with value zero. This puts $F$ in the centered exterior class.
	
	\subsection{Original roots and the actual boundary trace}
	\label{sec:smooth-exterior-potential}
	For $\rho>1$ the angular logarithmic potential
	$V_\rho(x)=(2\pi)^{-1}\int_0^{2\pi}
	\log|F(\rho e^{i\theta})-x|\,d\theta$ is
	\[
	V_\rho(x)=\frac1n\log\max(\rho^n,|Q(x)|).
	\]
	Indeed if $F(1/u)=1/u+g(u)$, the reciprocal function
	$1+u(g(u)-x)$ is analytic on the disk, has value $1$ at zero,
	and has no zero except a simple one at $1/w$ if $x=F(w)$.
	Apply Jensen's formula on radius $1/\rho$ and reverse the angular
	orientation. The logarithm of the product can be expanded off the
	single possible zero on the circle, an angular null set. The result
	is $\log\rho$ on $K$ and $\log\max(\rho,|w|)$ at $F(w)$.
	
	In particular, $V_\rho(x)\ge n^{-1}\log|Q(x)|$. Average over
	$x=s e^{i\alpha}$ for any $s>0$. The angular mean of
	$\log|y-s e^{i\alpha}|$ is $\log s+\log^+(|y|/s)$.
	Fubini applies to the double angular integral. The positive
	part is bounded on the compact product, and the slice integral has
	the lower bound $\log s$, which gives a uniform bound for the
	integral of the negative part. Factorization of $Q$ gives
	\[
	\frac1n\sum_i\log^+(|q_i|/s)
	\le \frac1{2\pi}\int_0^{2\pi}
	\log^+(|F(\rho e^{i\theta})|/s)\,d\theta.
	\]
	Here $\log^+0$ denotes zero. For $p>0$ and $a\ge0$,
	$a^p=p^2\int_{\mathbb R} e^{pt}\log^+(a e^{-t})\,dt$.
	For positive $a$ this follows by substituting $v=\log a-t$ and
	integrating $v e^{-pv}$ on $(0,\infty)$. For $a=0$ both sides
	vanish. Tonelli gives the $p$-moment comparison on the smooth
	circle. The actual radial trace and the bound
	$|F(\rho e^{i\theta})|\le2\rho$ justify dominated convergence as
	$\rho\downarrow1$. In particular
	$n^{-1}\sum_i|q_i|^p\le R_p(F)$ for every $p>0$.
	
	\subsection{Derivative roots, strictness, and sharpness}
	The de Bruijn--Springer comparison \cite[Theorem~4]{DBS1947}
	holds for $p\ge1$. We give a proof using the differentiator compression of \cite{Pereira}.
	For a degree $N$ polynomial with roots $\lambda_i$ listed with
	multiplicity, let $D=\operatorname{diag}(\lambda_i)$ and let
	$e=N^{-1/2}(1,\ldots,1)$. In an orthonormal basis with first vector
	$e$, compress $D$ to the last $N-1$ coordinates to obtain $A$.
	The determinant of $zI-A$ is the first diagonal cofactor of $zI-D$
	in this basis. Transforming the adjugate under change of basis
	gives the cofactor
	\[
	\frac1N\sum_i\prod_{\ell\ne i}(z-\lambda_\ell)
	=\frac1N\frac{d}{dz}\prod_i(z-\lambda_i).
	\]
	Thus the diagonal of a unitary Schur triangularization of $A$ lists
	the derivative roots $\mu_j$, with algebraic multiplicity. In
	ambient coordinates their diagonal formulas are
	$\mu_j=\sum_i|v_j(i)|^2\lambda_i$. The weights sum to one in each
	column. Across the $N-1$ columns the total in row $i$ is
	$1-|e(i)|^2=(N-1)/N$. The triangle inequality and Jensen for
	$r\mapsto r^p$ give
	$|\mu_j|^p\le\sum_i |v_j(i)|^2 |\lambda_i|^p$.
	Sum and divide by $N-1$, then iterate while the remaining degree
	is positive. This proves
	\[
	\frac{1}{n-j}\sum_{Q^{(j)}(z)=0}|z|^p
	\le \frac1n\sum_i|q_i|^p
	\le R_p(F) \le S(p).
	\]
	The last bound follows from Theorem~\ref{thm:radial-continuous}.
	Equality there would give $F(w)=w+\lambda/w$ with $|\lambda|=1$
	on $|w|>1$. Then the polynomial identity
	$w^nQ(w+\lambda/w)=w^{2n}$ is impossible. The left side has
	constant term $\lambda^n$ because $Q$ is monic and all lower
	terms give strictly positive powers of $w$. The right side has
	zero constant term. This proves strictness.
	
	For sharpness take $Q_n(z)=2T_n(z/2)$, where $T_n$ is the ordinary
	Chebyshev polynomial. It is monic and centered and its
	$2$-lemniscate contains the interval $[-2,2]$, which joins all its
	zeros. If the sublevel set split into two nonempty relatively
	open and closed sets, openness and properness of the restricted
	polynomial would make both images equal the whole value disk.
	Both sets would then meet the zero fiber, contradicting that the
	interval joins all of it. The standard cosine
	grid for the zeros of $T_n$ tends to the arcsine distribution.
	For fixed $j$ the zeros of $T_n^{(j)}$ have the same limit by repeated
	Rolle interlacing. The radius $R=2^{1/n}$ tends to one. The
	$p$-moment therefore converges to
	$\frac{1}{2\pi}\int_0^{2\pi}|2\cos\theta|^p\,d\theta=S(p)$.
	Together with strictness this proves the exact supremum and
	nonattainment.
	\section{Factor norms and regular stars}\label{sec:stars}
	For a nondegenerate plane continuum $K$ let $C(K)$ be the least
	constant for which
	\[
	\prod_{i=1}^r\|p_i\|_K\leq C(K)^{\sum_i\deg p_i}
	\Big\|\prod_{i=1}^rp_i\Big\|_K
	\]
	for every finite collection of nonzero polynomials.  Set
	$S_m=\bigcup_{j=0}^{m-1}e^{2\pi ij/m}[0,4^{1/m}]$;
	it has logarithmic capacity one.
	The next theorem answers the regular-star question posed after
	Theorem~2.3 of \cite{PR2} for plane continua.
	
	\begin{theorem}[Regular-star parity]\label{thm:parity}
		For every integer $m\geq2$, the comparison
		$C(K)\leq C(S_m)$ holds for every nondegenerate plane continuum
		with $m$-fold rotational symmetry if and only if $m$ is even.
		For $m=2k$, equality holds precisely for similarity images of
		$S_{2k}$. Translate the symmetry center to zero, fill the bounded
		complementary components, and scale to capacity one. If $b_{2k-1}$
		is the \emph{forward} coefficient of $z^{-(2k-1)}$ in the
		normalized exterior map of the filled set, then
		\begin{equation}\label{eq:star-quant}
			\log\frac{C(S_{2k})}{C(K)}
			\geq\frac{1-k^2|b_{2k-1}|^2}{96k^2}.
		\end{equation}
		For every odd $m\geq3$ there is an $m$-fold rotationally symmetric continuum
		$K_m$ with $C(K_m)>C(S_m)$.
	\end{theorem}
	\begin{proof}
		Translate the rotation center to zero, fill the bounded complementary
		components, and normalize the capacity to one as in
		Section~\ref{sec:exterior-normalization}. Filling preserves polynomial
		norms by the maximum modulus principle, and each similarity preserves
		the factor constant by polynomial composition. The filled set contains
		the rotation center. If equality forces the filled set to be a star,
		the star's empty interior gives equality with the original continuum
		in the same coordinates.
		The farthest-point formula \eqref{eq:actual-factor-formula} (see also \cite[Theorem 1.1]{PR2})
		reads
		\[
		\log C(K)=\int\log d_K(w)\,d\mu_K(w),\qquad
		d_K(w)=\max_{v\in K}|w-v|.
		\]
		Suppose $m=2k$.  Rotational covariance of the normalized exterior
		map gives $F(\eta z)=\eta F(z)$ for $\eta=e^{\pi i/k}$.  Put
		$E=\{w^k:w\in K\}$.  It is compact, connected, contains zero,
		has capacity one, and satisfies $E=-E$.  Thus every hypothesis of
		the cited symmetric comparison holds.  Moreover,
		$G(z^k)=F(z)^k$ is its normalized odd exterior map.  The power map
		pushes $\mu_K$ to $\mu_E$.  It sends $S_{2k}$ to $[-2,2]$.
		Writing $F(z)=z+b_{2k-1}z^{1-2k}+\cdots$ shows directly that
		\[
		G(\zeta)=\zeta+\beta\zeta^{-1}+\cdots,
		\qquad\beta=k b_{2k-1}.
		\]
		The omitted-point coefficient bound applied to $G$ gives
		$|v|\leq2$ for $v\in E$, hence $|w|\leq2^{1/k}$ for $w\in K$.
		Let $a=2^{1/k}$.  The function
		$\Phi(x)=\log(a+e^{x/k})-e^{2x}/(64k^2)$ is convex on
		$(-\infty,\log2]$: with $y=e^{x/k}/a\leq1$, the second derivative
		of the logarithm is
		$y/[k^2(1+y)^2]\geq y^{2k}/(4k^2)=e^{2x}/(16k^2)$.
		Extend $\Phi$ by its tangent at $\log2$ and apply the symmetric
		logarithmic comparison from Section~\ref{sec:verified-green-comparison}.
		This is the comparison of \cite[Corollary~6.3]{BLP} for the admissible
		test used here. As
		$d_K(w)\leq a+|w|$ and equality holds for the opposite tip on the
		star, it yields
		\[
		\log\frac{C(S_{2k})}{C(K)}\geq
		\frac{2-\int_E|v|^2\,d\mu_E(v)}{64k^2}.
		\]
		Oddness makes the remaining negative modes of $G$ start at index
		three.  Parseval and the area theorem therefore give
		$\int_E|v|^2\,d\mu_E(v)
		\leq1+|\beta|^2+(1-|\beta|^2)/3$.
		The corresponding actual odd exterior-map boundary estimate, including
		the passage to its radial trace, is kernel checked independently of
		the factor-norm formula and the all-order inverse-coefficient problem.
		Substitution proves \eqref{eq:star-quant}.  If $C(K)=C(S_{2k})$,
		then $|\beta|=1$; area rigidity gives $G(\zeta)=\zeta+\beta/\zeta$.
		Its omitted set is a rotated segment, whose full power preimage is
		a rotated $S_{2k}$.  The converse follows by similarity invariance.
		
		Now fix odd $m\geq3$, write $s=\sqrt{1-t^2}$, and for small real $t$
		define
		\[
		G_t(z)=\frac{(z-it+s)^2}{z-it},\qquad
		E_t=\C\setminus G_t(\{|z|>1\}),\qquad
		K_t=\{w:w^m\in E_t\}.
		\]
		If $G_t(z_1)=G_t(z_2)$ with $z_1\ne z_2$, the preimages are paired by
		$z\mapsto(1+itz)/(z-it)$, and
		$|z-it|^2-|1+itz|^2=(1-t^2)(|z|^2-1)$.
		Thus $G_t$ is exterior univalent.  The two critical points
		$it\pm s$ lie on the unit circle and have values $0,4s$.
		The pairing maps the circle to itself and identifies its two arcs
		between these points.  Each arc maps injectively except at its
		endpoints, so the omitted set $E_t$ is their common simple image arc.
		In particular it contains zero; its full $m$th-root preimage $K_t$
		is a continuum. The $m$ lifts of the arc with its zero endpoint
		removed are pairwise disjoint, and their closures meet only at zero.
		They have exactly $m$ outer endpoints. Any rotational symmetry
		preserves the unique branch point $0$ and permutes those endpoints.
		The normalized exterior map of $K_t$ is
		$F_t(z)=z\exp((\log(G_t(z^m)/z^m))/m)$. In the coordinate
		$u=1/z$ the argument of the logarithm is nonvanishing and analytic
		on $\D$ with value $1$ at zero. Choose the analytic logarithm with
		value zero at $u=0$. This gives $F_t(z)=z+O(z^{1-m})$. The identity $F_t(z)^m=G_t(z^m)$
		and the covariance $F_t(\eta z)=\eta F_t(z)$ for $\eta^m=1$
		show injectivity and surjectivity onto $\mathbb C\setminus K_t$.
		It follows that $K_t$ is a
		capacity-one continuum with all tips $R_t\eta$, $\eta^m=1$,
		$R_t=(4s)^{1/m}$, and $K_0=S_m$.
		With $\zeta=e^{i\theta}$ one has
		\[
		G_t((s-it)\zeta)=(2+2\cos\theta)L_t(\zeta),
		\qquad L_t(\zeta)=\frac{s-it}{1-it/((s-it)\zeta)}.
		\]
		As $t\to0$, $L_t\to1$ uniformly on the unit circle.
		For small $|t|$ take its logarithm in the right half-plane with
		$\log 1=0$, so $\log L_0=0$ at every angle.
		Set $r=(2+2\cos\theta)^{1/m}$,
		$Z_t=r\exp((\log L_t)/m)$ (zero at $\theta=\pi$), and
		$\delta=\pi-\pi/m$. The source-circle map $z\mapsto z^m$ followed
		by rotation by $(s-it)^{-1}$ preserves normalized angular measure.
		The radial boundary value of $F_t$ is a root of
		$G_t((s-it)e^{i\theta})$, so is a rotation of $Z_t$ by an $m$th
		root of unity. The farthest distance to $K_t$ is invariant under
		such rotations. Using the tip $R_te^{i\delta}$ gives
		\[
		\log C(K_t)\geq A(t)=\frac1{2\pi}\int_0^{2\pi}
		\log|Z_t-R_te^{i\delta}|\,d\theta.
		\]
		At $t=0$, $Z_0=r$ and there is equality. The squared distance on
		each radial segment is convex, and the tips at angles $\pm\delta$
		are farther than the other tips and the origin.
		Moreover $D=r^2+R_0^2-2rR_0\cos\delta\geq R_0^2$.
		Uniform continuity keeps the distance away from zero for small
		$|t|$, and the $t$-derivative of the integrand is bounded there.
		This justifies differentiation under the integral. Since $R'_0=0$ and
		$\partial_t\log L_t|_{t=0}=\sin\theta+i(\cos\theta-1)$,
		the radial term integrates to zero and
		\[
		A'(0)=\frac{R_0\sin\delta}{2\pi m}
		\int_0^{2\pi}\frac{r(1-\cos\theta)}{D}\,d\theta>0.
		\]
		Thus $C(K_t)>C(S_m)$ for sufficiently small positive $t$.
	\end{proof}
	
	
	\appendix
	\section{Exterior normalization and equilibrium estimates}
	\subsection{Exterior normalization of a plane continuum}
	\label{sec:exterior-normalization}
	
	Let \(L\) be a nondegenerate compact connected set after filling
	its bounded complementary components. For \(a\in L\), the reciprocal
	coordinate \(u=(z-a)^{-1}\) carries the exterior, including infinity,
	to an open connected set \(U\) containing zero. Its spherical
	complement is the image of \(L\) under the same M\"obius
	transformation, with \(a\) mapped to infinity. It is connected and
	contains at least one finite point.
	
	Here is a direct primitive argument for this domain. If \(P\subset U\)
	is compact, its filling is contained in \(U\). Indeed, the open
	unbounded component of \(\C\setminus P\), together with infinity, and
	the union of the bounded components separate the complement of \(P\)
	on the sphere. The connected spherical complement of \(U\) lies in
	the first set because it contains infinity.
	
	Any two points of an open connected plane domain can be joined by a
	finite horizontal and vertical walk. The points reachable from a
	fixed point form a relatively open and relatively closed subset,
	since a small rectangle contained in the domain can be crossed in two
	axis-parallel segments. Let \(P\) be the trace of a closed walk in
	\(U\). Subdivide its edges at the real and imaginary coordinates of
	all vertices. Add an outer row and column on each side of the
	resulting rectangular grid. Write \(H_{ij},V_{ij}\) for the signed
	counts of rightward and upward traversals. At each vertex the total
	incoming count equals the total outgoing count. Thus with
	\[
	W_{ij}=\sum_{t=0}^j H_{it}
	\]
	we have \(H_{ij}=W_{ij}-W_{i,j-1}\) and
	\(V_{ij}=W_{i-1,j}-W_{ij}\). These identities use zero weights outside
	the grid. The horizontal identity follows by subtraction and the
	vertical identity follows by summing the vertex balance up the
	column. Along the top row the vertical counts vanish, so the top
	weights coincide with the zero weight of the left outer column.
	
	The union of closed cells with \(W_{ij}\ne0\) is compact and has
	boundary contained in \(P\). At a point off \(P\), two incident
	cells sharing an edge have the same weight, because their shared
	edge has zero traversal count. The same holds for cells meeting at
	a vertex by passing through the other incident cells. Each such
	point is therefore interior to the union whenever it lies in it.
	Each complementary component of \(P\) meeting this union stays in
	the union and is bounded. All the cells with nonzero weight lie in
	the filling of \(P\) and hence in \(U\).
	
	An integral along the walk is now the sum of the rectangle integrals
	with weights \(W_{ij}\), by cancellation of the two contributions
	to each internal edge. The Cauchy theorem for rectangles makes this
	sum zero for a holomorphic integrand. Integrals from a base point
	are consequently independent of the walk. On a small disk they
	differ from the usual local primitive by a constant, which gives
	a global primitive. For a nonvanishing holomorphic \(q\), integrate
	\(q'/q\) and choose the constant of integration so that \(e^P=q\)
	at the base point. Since \((qe^{-P})'=0\), the identity holds on
	the domain and \(e^{P/2}\) is a square root.
	
	For completeness, this square-root property yields a disk
	uniformization as follows. For any omitted finite point \(a_0\),
	take a branch of \(\sqrt{z-a_0}\). It is injective and its open
	image is disjoint from the negative of the image. A small closed
	disk in the negative image is therefore omitted. Inversion at its
	center followed by scaling gives an embedding into the unit disk.
	A disk automorphism sends the image of the base point to zero.
	Maximize the derivative magnitude there over these normalized
	embeddings. Cauchy estimates and compact exhaustion give local
	uniform convergence of a maximizing subsequence. The limit is
	holomorphic, maps into the disk by the maximum principle, and has
	nonzero derivative at the base point. It is injective by the
	zero-free limit principle applied locally to differences of the
	embeddings. For that principle, a nonconstant limit with an
	isolated zero is bounded away from zero on a small surrounding
	circle. The maximum modulus principle bounds the reciprocals of
	the approximating functions on the disk, contradicting convergence
	to zero at the center.
	
	If the extremal embedding \(f\) omits \(b\) in the unit disk, put
	\(\varphi_b(w)=(w-b)/(1-\overline b w)\) and take
	\(g^2=\varphi_b\circ f\). The function \(g\) is an embedding into
	the disk. If \(d\) is its value at the base point,
	\(G=\varphi_d\circ g\) is normalized. Then \(f=\psi\circ G\) with
	\(\psi=\varphi_{-b}\circ(w\mapsto w^2)\circ\varphi_{-d}\).
	This disk self-map fixes zero and fails to be injective, so the
	strict Schwarz inequality gives \(|\psi'(0)|<1\). The chain rule
	contradicts extremality. The inverse of \(f\) uniformizes \(U\)
	by the disk and fixes zero when the base point is zero.
	
	Return to the filled continuum with an even rotational order and
	center zero. If zero were in the exterior, the rotation would
	induce a nontrivial disk automorphism with distinct fixed points
	corresponding to zero and infinity. Schwarz rules this out. We can
	therefore use the reciprocal coordinate at \(a=0\). Write its
	disk uniformization as \(u=f(t)\), with \(f(0)=0\). Conjugating
	central reflection by \(f\) gives a disk involution fixing zero.
	Its multiplier is \(-1\) by Schwarz and the fact that the
	involution is nontrivial. Thus \(f\) is odd. With
	\(a_1=f'(0)\ne0\), the formula
	\[
	F(w)=a_1/f(1/w)
	\]
	gives an exterior univalent map with expansion \(w+O(w^{-1})\).
	Its omitted set is \(a_1 L\). A rotation of the normalized map
	and its argument replaces \(a_1 L\) by \(|a_1|L\).
	Section~\ref{sec:radial-equilibrium} identifies the capacity of
	this normalized omitted set as one.
	\subsection{Logarithmic potential of the actual radial law}\label{sec:radial-potential}
	
	Let $F$ be centered exterior univalent, let $K=\mathbb C\setminus
	F(\{|z|>1\})$, and let $\nu_F$ be the image of normalized angular
	measure under its actual radial trace. For $0<s<1$, put
	\[
	K_s=\mathbb C\setminus F(\{|z|>s^{-1}\}),\qquad
	\nu_s=(\theta\mapsto F(s^{-1}e^{i\theta}))_\#\frac{d\theta}{2\pi}.
	\]
	The radial limit gives weak convergence of $\nu_s$ to $\nu_F$. The
	probability $\nu_F$ is carried by $K$. Indeed, if a radial limit were
	an exterior image point, the local inverse would force its source
	coordinate to be both in the open exterior and on the unit circle.
	
	Write $U_s(x)=\int\log|z-x|\,d\nu_s(z)$. Jensen's formula for the
	reciprocal germ $u(F(1/u)-x)$ gives $U_s(x)=-\log s$ on $K_s$.
	More generally $U_s(F(w))=\log\max(s^{-1},|w|)$ when $|w|>1$.
	At such an exterior image point the near-boundary circles are
	uniformly separated from the point. Dominated convergence gives
	$\int\log|z-F(w)|\,d\nu_F(z)=\log|w|$.
	
	For a point $x\in K$, fix $0<\varepsilon\le 1$ and use
	$\log\max(\varepsilon,|z-x|)$ as a test function. The supports are
	uniformly bounded near $s=1$, so we may cap the distance at $5$ to
	make the test bounded on the plane. Its $\nu_s$ integral is at
	least $-\log s$. Passing to $\nu_F$ shows that the truncated
	potential is nonnegative. For $\varepsilon=e^{-n}$, subtract
	the test from $\log 5$ and apply monotone convergence. The resulting
	nonnegative kernel is infinite at $z=x$, and its integral is at
	most $\log 5$. Thus $\nu_F\{x\}=0$, the logarithmic potential at
	$x$ is integrable, and its integral is nonnegative.
	
	To bound that integral from above, fix $0<s_0<1$. The point $x$ is
	interior to $K_{s_0}$. The images of bounded closed source annuli
	are compact, while $F(w)=w+o(1)$ at infinity. Together these facts
	separate $x$ from the image of $\{|w|\ge s_0^{-1}\}$. Choose
	$\varepsilon>0$ small enough that the circle of radius $\varepsilon$
	centered at $x$ lies in $K_{s_0}$. For $s_0<s<1$, the mean over
	that circle of $U_s$ is at most $-\log s_0$, by the two Jensen
	formulas above. The same mean equals
	$\int\log\max(\varepsilon,|z-x|)\,d\nu_s(z)$, by the elementary
	logarithmic circle average and Fubini's theorem. The logarithmic
	angular kernel is integrable on the product of the two circles.
	Its positive part is bounded on their compact images, and its
	negative part is controlled by the explicit logarithmic circle
	averages. Clipping at $5$ preserves the upper bound. Passing to
	$\nu_F$ and then letting $s_0$ tend to $1$ gives
	\begin{equation}\label{eq:radial-zero-potential}
		\int\log|z-x|\,d\nu_F(z)=0\qquad (x\in K).
	\end{equation}
	The same bounded-positive-part argument and \eqref{eq:radial-zero-potential}
	give an integrable mutual kernel with integral zero against
	$\lambda\otimes\nu_F$ for every probability $\lambda$ carried by $K$.
	The diagonal has zero mass for this product.
	\subsection{The energy minimizing property}\label{sec:radial-equilibrium}
	
	On $K\times K$, distances are at most $4$. For finite distances
	$0<d\le 4$, use the shifted logarithmic kernel
	\[
	W(d)=\log(25/d^2)=\int_0^\infty
	\frac{e^{-t d^2}-e^{-25t}}{t}\,dt.
	\]
	The identity is Frullani's integral. Its integrand is nonnegative
	and the integral is infinite at $d=0$, as follows by letting $d$
	decrease to zero. This shift keeps every energy integral in the
	extended nonnegative reals. Subtracting $2\log 5$ and dividing by
	$2$ recovers the usual negative logarithmic energy for probabilities.
	
	For probabilities $\alpha,\beta$ carried by $K$, put
	$H_t(\alpha,\beta)=\iint e^{-t|z-w|^2}\,d\alpha(z)d\beta(w)$.
	If $\gamma$ is the standard Gaussian law on $\mathbb R^2$, its
	characteristic function gives, for $t>0$,
	\begin{align*}
		&H_t(\alpha,\alpha)+H_t(\beta,\beta)
		-H_t(\alpha,\beta)-H_t(\beta,\alpha)\\
		&\qquad=\int_{\mathbb R^2}
		\left|\widehat\alpha(\sqrt{2t}\,v)-
		\widehat\beta(\sqrt{2t}\,v)\right|^2\,d\gamma(v)\ge0.
	\end{align*}
	The expectation of $\exp(i\sqrt{2t}\langle v,z-w\rangle)$ under
	$\gamma$ is $e^{-t|z-w|^2}$. Fubini's theorem for this bounded
	kernel gives the displayed identity. Its right-hand side vanishes
	only if the continuous characteristic functions agree everywhere,
	since every nonempty open set has positive Gaussian measure.
	The uniqueness theorem for characteristic functions then gives
	$\alpha=\beta$.
	
	Apply Tonelli's theorem to $W$. The mutual energy statement following
	\eqref{eq:radial-zero-potential} makes the two mixed $W$ energies
	of $\nu_F$ and any probability $\lambda$ on $K$ equal to $2\log 5$.
	The self $W$ energy of $\nu_F$ also equals $2\log 5$. Integrating the
	nonnegative Gaussian difference with weight $dt/t$ shows that the
	self $W$ energy of $\lambda$ is at least $2\log 5$, allowing the
	value $+\infty$. In the equality case all four energies are finite
	and the Gaussian difference vanishes for almost every positive $t$.
	Taking one such scale proves $\lambda=\nu_F$. Hence $\nu_F$ is the
	unique logarithmic equilibrium probability of the actual omitted set
	and $\operatorname{cap}(K)=1$.
	
	\subsection{The lower bound for admissible factor constants}
	
	Let $d_K(z)=\max_{w\in K}|z-w|$ and suppose $K$ has two distinct
	points. Fix $0<s<1$. For $N$ equally spaced angles $\theta_i$, set
	$a_i=F(s^{-1}e^{i\theta_i})$ and $Q_N(z)=\prod_i(z-a_i)$. These
	points lie outside $K$. On the compact set $K$ times the angle
	interval, the function
	\[
	(x,\theta)\longmapsto
	\log d_K(F(s^{-1}e^{i\theta}))-\log|F(s^{-1}e^{i\theta})-x|
	\]
	is continuous. Its equally spaced Riemann sums converge uniformly
	in $x$. The angular mean is independent of $x\in K$, by the smooth
	Jensen potential, and equals
	$\int\log d_K\,d\nu_s+\log s$. Evaluate the sum at a point of $K$
	maximizing $|Q_N|$. The defining factor inequality applied to the
	$N$ linear factors bounds that normalized sum by $\log c$ for
	each admissible constant $c$. Passing to the limit proves
	\[
	\int\log d_K\,d\nu_s+\log s\le \log c.
	\]
	The near-boundary outer laws are carried by a common compact disk,
	and $d_K$ is bounded below by half the distance between the two
	distinct points. Consequently bounded continuous tests give
	$\int\log d_K\,d\nu_F\le\log c$ as $s\uparrow1$.
	The reverse inequality follows from auxiliary polynomials and the
	Bernstein--Walsh bound.
	
	\subsection{The factor constant on an actual exterior omitted set}
	\label{sec:actual-factor-formula}
	
	Let $K$ be the nondegenerate omitted set of a normalized centered
	exterior map $F$, and let $\nu_F$ be its actual radial probability.
	The normalized farthest-point formula of Pritsker and Ruscheweyh
	\cite[Theorem~1.1]{PR2} is
	\begin{equation}\label{eq:actual-factor-formula}
		C(K)=\exp\left(\int\log d_K(z)\,d\nu_F(z)\right).
	\end{equation}
	The lower bound for every admissible constant was proved above.
	For the reverse bound, set
	$g_K(b)=\int\log|z-b|\,d\nu_F(z)$. In the exterior image this is
	$\log |F^{-1}(b)|$, and on $K$ it is zero. For a monic polynomial $P$
	of degree $n$, the maximum principle applied to
	$w^{-n}P(F(w))$ on $|w|\ge s^{-1}>1$, including infinity, gives
	\[
	\log|P(b)|\le\log\|P\|_K+n g_K(b).
	\]
	To see the limiting boundary bound, the curves $F(|w|=s^{-1})$ lie in
	a fixed compact disk for $s$ near one. Every cluster point as
	$s\uparrow1$ is in $K$, since an exterior limit would have an inverse
	of modulus greater than one. The displayed inequality is also valid
	at zeros with the convention $\log0=-\infty$.
	
	Write $P(z)=\prod_{i=1}^n(z-a_i)$. For an auxiliary polynomial $Q$
	of degree $m$, with nonzero leading coefficient and $Q(a_i)\ne0$,
	apply the preceding estimate to $P$ at the roots of $Q$, counting
	multiplicities. The leading coefficient of $Q$ cancels after writing
	the pairwise distances in the two orders. This gives
	\begin{equation}\label{eq:factor-auxiliary-polynomial}
		\sum_{i=1}^n\log|Q(a_i)|\le m\log\|P\|_K
		+n\int\log|Q(z)|\,d\nu_F(z).
	\end{equation}
	Each integral is finite by the pointwise logarithmic integrability of
	$\nu_F$ established in \S\ref{sec:radial-equilibrium} and its
	preceding potential calculation.
	
	Choose a finite $\delta$-net $b_1,\ldots,b_N$ in $K$, including two
	distinct points, and let $Q(z)=\sum_{j=1}^N\xi_j(z-b_j)^m$, where
	the $\xi_j$ are independent standard complex Gaussians with real and
	imaginary variances one. The leading coefficient is nonzero almost
	surely. For fixed $z$, $Q(z)$ has the law of a standard complex
	Gaussian multiplied by
	$\sigma(z)=(\sum_j |z-b_j|^{2m})^{1/2}>0$. Thus the expectation of
	$\log|Q(z)|$ is $\log\sigma(z)+c_0$, with the same $c_0$ at every $z$.
	The expectation of its absolute value is at most
	$|\log\sigma(z)|+\mathbb E|\log|\xi_1||$. The last expectation is
	finite by the Gaussian density in polar coordinates, and
	$\log\sigma$ is bounded on $K$. Fubini therefore applies to the
	right side of \eqref{eq:factor-auxiliary-polynomial}. The two Gaussian
	constants cancel after taking expectations. Dividing by $m$ yields
	the same inequality with $\log|Q|$ replaced by
	$m^{-1}\log\sigma$ and the coefficient of $\log\|P\|_K$ reduced to one.
	
	If $d_B(z)=\max_j|z-b_j|$, then $d_B\le d_K\le d_B+\delta$ and
	\[
	\log d_B(z)\le m^{-1}\log\sigma(z)\le
	\log d_B(z)+\frac{\log N}{2m}.
	\]
	Half the distance between the two fixed points of the net bounds
	$d_B$ below everywhere. First taking $\delta$ small with those two
	points fixed, then taking $m$ large, gives arbitrarily small uniform
	error relative to $\log d_K$. Exponentiation proves the factor
	inequality for all finite linear families with the constant in
	\eqref{eq:actual-factor-formula}. Factoring each nonzero polynomial
	over $\mathbb C$ and applying submultiplicativity of the uniform norm
	proves admissibility for arbitrary families. A zero factor makes both
	sides zero. This proves \eqref{eq:actual-factor-formula}.
	\subsection{A direct Green comparison for the even power quotient}
	\label{sec:verified-green-comparison}
	
	For an actual odd normalized exterior map $G$ with connected omitted
	set $E$, write $g(z)=\log|G^{-1}(z)|$ on the exterior image and $g=0$
	on $E$. Oddness makes $g(-z)=g(z)$ and places $0$ in $E$: a preimage of
	zero in the open exterior disk would have a distinct negative
	preimage. The zero extension is continuous by the outer collars and
	is harmonic wherever it is positive. Let $g_L$ be the corresponding
	Green function for the segment $L=[-2,2]$. We show that the circular
	mean of $g$ on every centered circle is at most that of $g_L$.
	This is the symmetric case of the comparison used in
	\cite[Section~6]{BLP}.
	
	Work in the logarithmic chart $z\mapsto i\exp(z/2)$ and set
	$u_\tau(x,t)=(g(i\exp((x+it)/2))-\tau)_+$ for $\tau>0$.
	For $0\le h\le\pi$, let $U_\tau(x,h)$ be the supremum of the
	integral of $u_\tau(x,\cdot)$ over measurable angular sets of length
	$2h$ in one turn. This supremum is independent of the turn by
	periodicity. It is attained by taking the strict superlevel at a
	quantile and a subset of its level of the required remaining length.
	Equivalently, it is the infimum over real $a$ of
	\[
	2ha+\int_{-\pi}^{\pi}(u_\tau(x,t)-a)_+\,dt.
	\]
	The two formulas give continuity in $(x,h)$ up to the angular edges.
	For instance, a uniform change of size $\delta$ in the angular
	profile changes a top integral by at most $2\pi\delta$, and changing
	the length changes the integral by at most the length difference
	times a common uniform bound for the profiles.
	
	For $0<h<\pi$ one can choose a maximizing set containing a short
	interval and avoiding margins around the angular cut. Indeed positive
	levels of the untruncated profile are null unless that angular
	profile is constant. The profile is real analytic at positive
	values, and a positive level that is locally constant propagates
	through its positive component by analyticity and continuity. A
	positive-level quantile therefore leaves a strict lower interval
	for the cut and includes a strict upper interval. At a zero-level
	quantile of $u_\tau$, the set where $g(i\exp((x+it)/2))<\tau$ has
	positive length. Choose the cut there and take the missing mass from
	the zero plateau, keeping short margins at the cut. A constant
	clipped profile admits an interval of the prescribed length.
	
	These sets give the local submean for $U_\tau$. For a small circular
	increment $(\rho\cos\theta,\pm\rho\sin\theta)$ with
	$0\le\theta\le\pi$, put $d=\rho\sin\theta$ and translate the
	maximizing set $A$ to $A-d$ and $A+d$. Split each translated set at a
	point $p$ inside the included interval. Join the left piece of $A-d$
	to the right piece of $A+d$, and join the two remaining pieces.
	Since $[p-d,p+d]\subset A$, the new masses are $2(h+d)$ and $2(h-d)$.
	Their combined indicators equal those of the two translates. The
	margins keep both candidates in one turn. At the original vertical
	centers where $u_\tau>0$, the
	unclipped profile is harmonic in a common neighborhood, since the
	positive cutoff keeps those centers in a compact subset of the
	positive domain. Thus $u_\tau$, the maximum with zero, satisfies the
	circle submean there. At a zero center the submean is immediate from
	nonnegativity. Integrating over the maximizing set and pairing
	opposite angular increments proves the claimed submean.
	
	Let $V(x,h)$ be the analogous top integral of the unclipped slit
	profile. The confocal ellipse formula
	\[
	\cosh g_L(z)=\frac{|z-2|+|z+2|}{4}
	\]
	shows that its angular top sets are the centered arcs $[-h,h]$ in
	this chart. On the open upper strip the arc integral is harmonic.
	In fact the holomorphic function
	$\log(-iG_L^{-1}(i\exp(z/2)))$ on $|\Im z|<\pi$ has real part
	equal to the slit Green profile and is real on the center line.
	The argument of this logarithm has positive real part. Integrating
	its real part on a vertical half-arc gives the difference of the
	imaginary parts of a local holomorphic primitive at the endpoints.
	The primitive's imaginary part is constant on the center line.
	
	At $h=0$ both top integrals vanish. At $h=\pi$, if the physical
	radius is at least the maximum radius of $E$, the full mean of $g$
	equals the logarithm of that radius. The full slit mean is at least
	this logarithm by Jensen's formula for the circular log kernel.
	Otherwise the circle meets $E$, by connectedness and $0\in E$.
	The clipped profile vanishes on a short arc there, so its full top
	integral stays constant on moving a short distance inward from
	$h=\pi$. The slit top integral is nondecreasing in $h$.
	As $x\to-\infty$ the profiles tend uniformly to zero. As
	$x\to+\infty$ each unclipped Green profile differs uniformly from
	$x/2$ by a quantity tending to zero, since the normalized exterior
	maps are $w+O(1/w)$ at infinity.
	
	On a wide compact strip rectangle compare $U_\tau-V$ after adding
	$-\epsilon h+\epsilon h^2/(4\pi)$. The quadratic gives a strict
	submean at interior points. At a top-edge point controlled by a
	zero arc, moving inward increases the perturbed difference. All
	other boundary values are at most $\epsilon$ once the rectangle
	is wide enough. The maximum principle gives
	$U_\tau-V\le\epsilon(1+\pi)$. Let $\epsilon\downarrow0$ and then
	$\tau\downarrow0$, using the uniform bound $2\pi\tau$ for the change
	in a top integral. At $h=\pi$ the chart covers a half-turn of the
	physical circle with twice the angular parameter. Odd symmetry gives
	the desired full-circle comparison.
	
	To pass from this comparison to the test in the even-order proof,
	let $\mu_E$ be the equilibrium probability measure. Capacity one gives
	$g(z)=\int\log|z-v|\,d\mu_E(v)$, and $g(0)=0$ gives
	$\int\log|v|\,d\mu_E(v)=0$. The same holds for $L$. The circular
	average of $\log|z-v|$ on $|z|=r>0$ equals
	$\max\{\log r,\log|v|\}$. Averaging the potential therefore gives,
	for every $a=\log r$, a comparison of the integrals of
	$(\log|v|-a)_+$ for $E$ and $L$. These tests, together with the
	equal means, give the convex comparison for the logarithmic radii
	by piecewise linear approximation. The test used above is convex
	on $(-\infty,\log 2]$; its tangent line continues it to the right.
	\section*{Verification}
	The proofs use the stated analytic comparisons and the displayed
	scalar estimates. The single multiplier reduces the energy bound
	to two explicit quadratic bounds. The final phase budgets use
	the rational bounds and elementary series estimates in
	Section~\ref{sec:large}. These calculations and the finite
	low-order coefficient expansions can be checked from the
	formulas in the text. The main statements also have an independent
	Lean~4 verification.
	
	\section*{Declaration of AI use}
	OpenAI's ChatGPT and Codex assisted with proof exploration, Lean~4
	formalization, and revision of the exposition.
	
	\begingroup\small
	\begin{thebibliography}{99}
		\bibitem{Springer} G. Springer,
		The coefficient problem for schlicht mappings of the exterior of the unit circle,
		\emph{Trans. Amer. Math. Soc.} \textbf{70} (1951), 421--450.
		\href{https://doi.org/10.1090/S0002-9947-1951-0041935-5}{doi:10.1090/S0002-9947-1951-0041935-5}.
		\bibitem{Kubota} Y. Kubota, Coefficients of meromorphic univalent functions,
		\emph{K\=odai Math. Sem. Rep.} \textbf{28} (1977), 253--261.
		\bibitem{Schober77} G. Schober, Coefficients of inverses of meromorphic
		univalent functions, \emph{Proc. Amer. Math. Soc.} \textbf{67} (1977), 111--116.
		\bibitem{Schober} G. Schober,
		Coefficients of inverses of univalent functions with quasiconformal extensions,
		\emph{Kodai Math. J.} \textbf{2} (1979), 411--419.
		\bibitem{BLP} A. Baernstein II, R. S. Laugesen, and I. E. Pritsker,
		Moment inequalities for equilibrium measures in the plane,
		\emph{Pure Appl. Math. Q.} \textbf{7} (2011), 47--82.
		\href{https://doi.org/10.4310/PAMQ.2011.v7.n1.a4}{doi:10.4310/PAMQ.2011.v7.n1.a4}.
		\bibitem{Circle} A. Baernstein II, Convolution and rearrangement on circle,
		\emph{Complex Variables Theory Appl.} \textbf{12} (1989), 33--37.
		Correction, \emph{ibid.} \textbf{26} (1995), 381--382.
		\bibitem{PR2} I. E. Pritsker and S. Ruscheweyh,
		Inequalities for products of polynomials II,
		\emph{Aequationes Math.} \textbf{77} (2009), 119--132.
		\bibitem{Pom} Ch. Pommerenke, \emph{Univalent Functions},
		Vandenhoeck \& Ruprecht, G\"ottingen, 1975.
		\bibitem{Pereira} R. Pereira, Differentiators and the geometry of
		polynomials, \emph{J. Math. Anal. Appl.} \textbf{285} (2003), 336--348.
		\href{https://doi.org/10.1016/S0022-247X(03)00465-7}{doi:10.1016/S0022-247X(03)00465-7}.
		\bibitem{DBS1947} N. G. de Bruijn and T. A. Springer,
		On the zeros of a polynomial and of its derivative II,
		\emph{Proc. K. Ned. Akad. Wet. Amsterdam} \textbf{50} (1947), 264--270,
		Theorem~4.
	\end{thebibliography}
	\endgroup
	
\end{document}